% !TEX root = guida_tecnica_firmware_tank_esp32.tex
\documentclass[11pt,a4paper]{article}

\usepackage[utf8]{inputenc}
\usepackage[T1]{fontenc}
\usepackage[italian]{babel}
\usepackage[a4paper,margin=16mm,headheight=14pt]{geometry}
\usepackage{array}
\usepackage{booktabs}
\usepackage{enumitem}
\usepackage{fancyhdr}
\usepackage{hyperref}
\usepackage{lastpage}
\usepackage{longtable}
\usepackage{listings}
\usepackage{xcolor}
\usepackage{titlesec}
\usepackage{url}

\definecolor{navy}{HTML}{10243F}
\definecolor{bluegray}{HTML}{EAF0F6}
\definecolor{codebg}{HTML}{F7F9FC}
\definecolor{codeframe}{HTML}{C8D2DD}
\hypersetup{colorlinks=true,linkcolor=navy,urlcolor=navy,pdftitle={Guida tecnica firmware Tank e ESP32}}
\setlist[itemize]{leftmargin=*,itemsep=3pt}
\setlist[enumerate]{leftmargin=*,itemsep=3pt}
\titleformat{\section}{\Large\bfseries\color{navy}}{\thesection}{0.7em}{}
\titleformat{\subsection}{\large\bfseries\color{navy}}{\thesubsection}{0.6em}{}
\pagestyle{fancy}
\fancyhf{}
\lhead{Guida tecnica Tank + ESP32}
\rhead{Snapshot firmware 19/07/2026}
\cfoot{Pagina \thepage\ di \pageref{LastPage}}
\fancypagestyle{plain}{
  \fancyhf{}
  \lhead{Guida tecnica Tank + ESP32}
  \rhead{Snapshot firmware 19/07/2026}
  \cfoot{Pagina \thepage\ di \pageref{LastPage}}
}
\setlength{\emergencystretch}{3em}
\newcommand{\codepath}[1]{\path{#1}}
\lstset{language=C++,basicstyle=\ttfamily\scriptsize,numbers=left,numberstyle=\tiny\color{navy},numbersep=7pt,breaklines=true,breakatwhitespace=false,columns=fullflexible,frame=single,framesep=5pt,rulecolor=\color{codeframe},backgroundcolor=\color{codebg},showstringspaces=false,tabsize=4}

\begin{document}
\begin{titlepage}
\centering
\vspace*{16mm}
{\Huge\bfseries\color{navy} Guida tecnica completa\\[3mm] Firmware Tank + controller ESP32\par}
\vspace{8mm}
{\large Lettura per blocchi, baseline e aggiornamento post-correzioni\par}
\vfill
\begin{tabular}{rl}
\textbf{Baseline della guida:} & \codepath{7f43f59c00b3c0022841e7f037ebe5a3e8a88235} \\
\textbf{Snapshot analizzata:} & working tree post-correzioni e refactor cingoli, 19 luglio 2026 \\
\textbf{Firmware coperto:} & \codepath{tank/} e \codepath{controller-esp32/} \\
\textbf{Sorgenti esclusi:} & CAD, librerie vendor e file README generici \\
\textbf{Documento collegato:} & \codepath{audit_firmware_tank_esp32.pdf} \\
\end{tabular}
\vfill
\fcolorbox{navy}{bluegray}{\parbox{0.88\textwidth}{\small
I listing vengono caricati direttamente dai sorgenti e quindi mostrano il firmware corrente alla compilazione del PDF. Le sezioni dei cingoli e della lettura joystick sono state riallineate riga per riga alla snapshot del 19 luglio 2026; per gli altri moduli il capitolo ``Aggiornamento post-correzioni'' descrive il delta rispetto alla baseline. I comportamenti hardware restano distinti da quelli dimostrabili dal software.}}
\end{titlepage}

\tableofcontents
\newpage

\section{Come leggere la guida}

Ogni file e' presentato con quattro livelli complementari: una panoramica, la tabella dei blocchi logici, un listing completo con numeri di riga e una tabella riga per riga. Il listing e' la fonte principale per vedere sintassi e commenti. Per 
\codepath{trackController.*} e \codepath{joystickReader.*}, tabelle e listing sono allineati alla snapshot corrente; per gli altri moduli modificati usare prima il capitolo ``Aggiornamento post-correzioni'', poi il listing corrente e infine la spiegazione baseline delle righe non cambiate.

I riferimenti \codepath{A-01} fino a \codepath{A-14} rimandano agli ID dell'audit allegato. Lo stato aggiornato (risolto, mitigato o aperto) e' nella tabella post-correzioni dell'audit; un riferimento non significa che la singola riga sia ancora errata.

\section{Architettura del sistema}

\begin{verbatim}
ESP32 controller
  ADC joystick + pulsanti
  -> WiFi station, UDP V1 / porta 4210
  -> Uno R4 WiFi: AP Tank_AP + UdpReceiver_update()
  -> TankCommand
     -> TrackController -> MotorController -> PWMController -> I2C/PCA9685 -> M3/M1
     -> ServoTorretta -> MotorController -> PCA9685 -> S5/S6
     -> StepperTorretta -> D2-D5 -> driver stepper
     -> Relay Fire -> D7
\end{verbatim}

\begin{longtable}{p{0.24\textwidth}p{0.26\textwidth}p{0.40\textwidth}}
\toprule
\textbf{Segnale} & \textbf{Origine/destinazione} & \textbf{Significato} \\
\midrule
\endhead
\codepath{driveX} & GPIO32 dopo lo swap -> tank & Sterzata differenziale. \\
\codepath{driveY} & GPIO34 dopo lo swap -> tank & Avanti/indietro. \\
\codepath{turretX} & GPIO33 dopo lo swap -> D2--D5 & Rotazione orizzontale. \\
\codepath{elevationY} & GPIO35 dopo lo swap -> S5/S6 & Elevazione servo. \\
\codepath{zero}, \codepath{fire} & GPIO25/GPIO26, LOW attivo & Azzeramento e impulso relay. \\
I2C PCA9685 & Uno R4 -> shield & PWM motori M1/M3 e servo S5/S6. \\
\bottomrule
\end{longtable}

\section{Tecnologie e concetti trasversali}

\subsection{PlatformIO, Arduino e C++}
PlatformIO legge \codepath{platformio.ini}, risolve framework/librerie e compila. Il modello Arduino chiama \codepath{setup()} una volta e poi \codepath{loop()} ripetutamente. Gli header \codepath{.h} dichiarano tipi e API; i \codepath{.cpp} definiscono il comportamento. Le include guard impediscono doppie dichiarazioni. Macro \texttt{\#define} sono sostituzioni del preprocessore: non hanno tipo e vanno usate con attenzione.

\subsection{WiFi AP, station e UDP}
Il tank crea l'access point; l'ESP32 e' una station. UDP invia datagrammi a una porta senza handshake applicativo: non garantisce consegna, ordine, unicita', autenticita' o conferma. Il frame V1 e' \codepath{V1;driveX;driveY;turretX;elevationY;zero;fire}. Un frame normale massimo e' 26 byte; il buffer TX da 64 byte e' adeguato. Su richiesta, il receiver conserva il parsing compatibile con la baseline: legge al massimo 63 byte, usa \codepath{sscanf} per sei interi e limita gli assi con \codepath{constrain}. Non rifiuta rigorosamente suffissi, campi extra o conversioni \codepath{int} fuori range; V1 inoltre non autentica mittente, sequenza, replay o freschezza. \codepath{WL_CONNECTED} e \codepath{udpReady} provano solo lo stato locale della station/socket, non che il tank abbia applicato il comando.

\subsection{ADC, INPUT\_PULLUP e joystick}
L'ESP32 legge tensioni analogiche con ADC. La risoluzione 12 bit produce normalmente 0--4095, poi \codepath{map} le adatta a 0--1023. \codepath{constrain} limita il valore. GPIO32--35 sono ADC1: questa scelta evita il conflitto WiFi tipico di ADC2. Con \codepath{INPUT_PULLUP}, il pin e' HIGH a riposo e LOW alla pressione. Il controller applica una deadzone di ingresso da 20 alla guida e da 80 alla torretta: filtrano il rumore del potenziometro attorno al centro calibrato. Il tank conserva una deadzone distinta da 100 sul comando cingolo, prima e dopo il mixing, per proteggere anche da sender UDP diversi dall'ESP32. Il firmware valida centro/spread alla calibrazione, applica debounce di 40 ms e richiede neutralita' prima dell'armamento; non puo' pero' diagnosticare ogni guasto di potenziometro o cablaggio dopo il boot.

\subsection{I2C, PCA9685 e PWM}
I2C e' un bus indirizzato: il master invia indirizzo, registro e byte, mentre un ACK indica soltanto che un dispositivo ha risposto. La PCA9685 offre 16 canali PWM a 12 bit; i registri ON/OFF descrivono il ciclo. Il valore 4096 viene usato come flag full-on/full-off. Tutti i canali condividono la stessa frequenza: qui 50 Hz e' corretto per servo ma un compromesso per motori DC. Il firmware ora propaga NACK, letture corte e canali non validi in un fault latched; una chiamata I2C gia' bloccata e l'ultimo PWM fisico richiedono comunque un enable hardware indipendente.

\subsection{Motori, servo, stepper e relay}
Un H-bridge decide verso e PWM dei motori DC. Il firmware ora porta PWM a zero prima di cambiare verso e il controller cingoli introduce un dead-time di 30 ms; non e' una rampa completa. Il servo interpreta la larghezza di un impulso periodico; la coppia S5/S6 usa un angolo specchiato. Lo stepper e' pilotato da quattro fasi D2--D5, senza feedback di posizione. Il relay traduce D7 in contatto elettrico: durata di impulso, polarita' e stato al reset sono requisiti safety e non solo dettagli software.

\subsection{millis, polling e interrupt}
Il progetto non definisce ISR o \codepath{attachInterrupt}. Tutto avviene per polling nel loop. Il pattern \codepath{millis() - last >= interval} e' corretto attraverso il rollover unsigned, ma una scadenza e' verificata solo se il loop torna a eseguire. WiFi, I2C, Serial e \codepath{delay} possono introdurre jitter. Nessuna ISR applicativa significa nessuna race \codepath{volatile} nel codice visibile, non significa real-time garantito.

\section{Contratto UDP e verifica incrociata}

\begin{longtable}{p{0.17\textwidth}p{0.30\textwidth}p{0.39\textwidth}}
\toprule
\textbf{Campo} & \textbf{Controller} & \textbf{Tank} \\
\midrule
\endhead
\codepath{driveX} & GPIO32 dopo swap, 0--1023 & Sterzo differenziale. \\
\codepath{driveY} & GPIO34 dopo swap, 0--1023 & Avanti/indietro. \\
\codepath{turretX} & GPIO33 dopo swap, 0--1023 & Stepper torretta. \\
\codepath{elevationY} & GPIO35 dopo swap, 0--1023 & Servo elevazione. \\
\codepath{zero} & GPIO25, LOW = 1 & Fronte Zero. \\
\codepath{fire} & GPIO26, LOW = 1 & Fronte relay Fire. \\
\bottomrule
\end{longtable}

Il sender serializza i sei campi nell'ordine mostrato e il receiver usa lo stesso ordine. Questa coerenza e' un aspetto positivo. A-08 resta aperto per scelta di compatibilita' del parser \codepath{sscanf}; restano inoltre A-04 e A-09, cioe' autenticazione, anti-replay, ordine, freschezza e conferma applicativa.

\section{Aggiornamento post-correzioni: cosa e' stato modificato}

\subsection{Come usare questo capitolo}

Questa e' la revisione della guida richiesta dopo l'implementazione dei fix. Il protocollo UDP V1, la porta 4210 e i pin fisici non sono cambiati. Sono invece cambiati undici moduli del firmware e alcune API locali: \codepath{PWMController::begin} e \codepath{MotorController::begin} ora restituiscono un esito, \codepath{MotorController} espone lo stato I2C e \codepath{JoystickReader} espone il riarmo neutro.

I listing sotto ``Lettura completa dei file'' usano \codepath{\textbackslash lstinputlisting} e sono quindi aggiornati automaticamente al sorgente corrente quando questo PDF viene compilato. Le tabelle riga per riga della baseline restano utili per gli elementi invariati; per ogni modulo nella tabella seguente, questa sezione prevale su vecchi numeri di riga, vecchie panoramiche e frasi che descrivono il comportamento prima dei fix.

\subsection{Mappa delle modifiche}

\begin{longtable}{p{0.27\textwidth}p{0.29\textwidth}p{0.30\textwidth}}
\toprule
\textbf{File} & \textbf{Cosa e' stato cambiato} & \textbf{Effetto e limite} \\
\midrule
\endhead
\codepath{tank/src/udpReceiver.cpp} & Stato \codepath{NETWORK_FAULT/NETWORK_READY}, retry AP/UDP, parser V1 \codepath{sscanf} compatibile e riarmo neutro. & Non esiste piu' \codepath{while(true)}; il parser accetta sei interi e clampa gli assi, quindi A-08 resta aperto. Una chiamata WiFiS3 puo' comunque bloccarsi internamente. \\
\codepath{tank/src/udpReceiver.h} & \codepath{TankCommand.connected} documenta ora ``comando autorizzato dal riarmo'', non soltanto traffico recente. & Il main deve trattare \codepath{false} come divieto assoluto di comandare attuatori. \\
\codepath{tank/src/main.cpp} & GPIO sicuri prima di Serial, guardia globale su link/shield, Fire e Zero con rilascio obbligatorio, fault I2C latched. & Spegne i segnali controllabili dal firmware; non sostituisce un enable elettrico esterno. \\
\codepath{tank/lib/PWMController/*} & Le primitive I2C restituiscono \codepath{bool} e memorizzano \codepath{communicationHealthy}. & Rileva NACK, letture corte e canali invalidi quando la chiamata ritorna; non sblocca Wire. \\
\codepath{tank/lib/motorController/*} & Inizializzazione e health I2C propagati; PWM zero prima del cambio direzione. & Evita la commutazione software con PWM attivo, ma non e' una rampa completa. \\
\codepath{tank/src/trackController.*} & Deadzone di sicurezza \codepath{TRACK_COMMAND_DEADZONE=100} prima/dopo mixing, stato M1/M3 e dead-time di 30 ms durante inversione. & Lo stato conserva verso e istante dell'ultimo stop anche dopo il ritorno al centro: una richiesta immediata del verso opposto non salta la pausa. Valori e frenatura fisica restano da misurare sul veicolo. \\
\codepath{controller-esp32/src/joystickReader.*} & Deadzone di ingresso separate (guida 20, torretta 80), calibrazione validata, debounce e riarmo neutro di 400 ms. & Le inversioni di Joy1 sono applicate solo qui; se la calibrazione fallisce il controller invia solo dati sicuri fino al riavvio. \\
\codepath{controller-esp32/src/udpSender.cpp} & Stato \codepath{udpReady}, controllo delle API UDP e soppressione del campione pre-reconnect. & Conferma soltanto lo stato locale del socket, non la ricezione/applicazione sul tank. \\
\bottomrule
\end{longtable}

\subsection{Flusso di sicurezza attuale}

\begin{verbatim}
ESP32: CALIBRAZIONE_VALIDA -> WAIT_NEUTRAL_400MS -> invio V1
                                             |            |
reconnect -----------------------------------+            v
Tank: NETWORK_FAULT -> retry AP/UDP -> WAIT_3_NEUTRI -> connected=true
  |                         |                         |
  +--> TankCommand sicuro <-+-------------------------+
                                                    -> attuatori
I2C NACK/lettura corta -> SHIELD_FAULT_LATCHED -> relay OFF, stepper OFF,
                                                   nessun nuovo comando PWM
\end{verbatim}

Il doppio riarmo e' intenzionale: il controller non trasmette subito un comando attivo dopo boot/reconnect e il tank non si fida comunque del primo datagramma che \codepath{sscanf} riesce a interpretare. Nessuna freccia del diagramma dimostra il taglio fisico dell'alimentazione: quello richiede hardware dedicato.

\subsection{Receiver UDP del tank: delta tecnico}

\begin{itemize}
\item \codepath{UdpReceiver_begin()} non avvia piu' AP e socket in modo definitivo nel \codepath{setup}. Azzera lo stato, lascia il comando sicuro e programma il primo tentativo dal \codepath{loop}.
\item \codepath{ensureNetworkReady()} tenta \codepath{WiFi.beginAP()} e \codepath{udp.begin()} al massimo una volta ogni 3 s dopo un errore restituito. Il timestamp viene aggiornato dopo la chiamata, per evitare retry immediati se la libreria ha atteso a lungo.
\item Il parser e' stato riportato alla semantica originale: \codepath{sscanf("V1;\%d;...", ...)} deve convertire sei interi; gli assi sono poi limitati con \codepath{constrain(..., 0, 1023)} e i pulsanti valgono true solo se uguali a 1.
\item La lettura copia al massimo 63 byte e aggiunge il NUL, ma non rifiuta esplicitamente payload oversize, suffissi, campi extra o conversioni \codepath{int} fuori range. Se \codepath{sscanf} non converte sei campi il watchdog non viene aggiornato; un datagramma troncato che contiene gia' sei conversioni puo' invece essere accettato.
\item \codepath{controlsArmed} diventa vero solo dopo tre frame validi con assi entro \(512 \pm 20\) e Fire/Zero rilasciati. Prima di allora \codepath{connected} resta falso e il main ferma tutto.
\end{itemize}

Resta un datagramma elaborato per chiamata; V1 non contiene sequence number, lease, HMAC o anti-replay. In particolare, ``retry controllato'' non vuol dire che \codepath{WiFi.beginAP()} o \codepath{udp.parsePacket()} siano non bloccanti.

\subsection{Main tank, Fire/Zero e fault PCA}

\begin{itemize}
\item Il latch inattivo di D7 e i latch LOW di D2--D5 vengono scritti prima di \codepath{Serial.begin()} e \codepath{delay(1000)}. Questo riduce la finestra software di boot, non garantisce il livello prima dell'esecuzione del firmware.
\item Se \codepath{command.connected} e' falso oppure \codepath{shieldReady} e' falso, il loop disabilita il relay e ferma i cingoli/stepper prima di elaborare movimento, Fire o Zero.
\item Fire e Zero usano lo schema ``wait release'': dopo fault/reconnect una pressione gia' mantenuta viene memorizzata ma non genera un fronte. Solo rilascio seguito da una nuova pressione e' operativo.
\item \codepath{enterShieldFault()} e' latched: dopo NACK o lettura I2C corta il firmware inibisce relay, stepper e nuovi comandi. Non tenta un riarmo automatico della shield.
\end{itemize}

\subsection{PCA, motori e inversione}

\begin{itemize}
\item \codepath{PWMController} controlla ora \codepath{Wire.endTransmission()} e \codepath{requestFrom()}; un errore rende \codepath{isCommunicationHealthy()} falso. \codepath{setPWM()} rifiuta inoltre canali oltre 15.
\item \codepath{MotorController::DCrun()} azzera PWM se il verso registrato cambia, poi scrive direzione e PWM. \codepath{DCbrake()} azzera anche il verso registrato.
\item \codepath{TrackController} mantiene per ogni cingolo l'ultimo verso richiesto, il verso prima dello stop, comando pendente e istante di stop. In inversione scrive brake, attende 30 ms e applica l'ultimo comando pendente; uno stop annulla l'inversione pendente ma conserva verso e tempo per non bypassare la pausa nel percorso centro -> verso opposto.
\end{itemize}

Questo limita il colpo software di inversione ma non prova che \codepath{DCbrake()} corrisponda a frenatura fisica, non crea una rampa di velocita' e non risolve il PWM DC condiviso a 50 Hz.

\subsection{Controller: calibrazione, pulsanti e socket UDP}

\begin{itemize}
\item La calibrazione usa media, minimo e massimo di 80 campioni. Ogni asse e' valido solo con centro entro \(512 \pm 70\) e spread non oltre 40; altrimenti \codepath{safeControllerData()} invia assi a 512 e pulsanti falsi fino al riavvio.
\item \codepath{DRIVE_INPUT_DEADZONE=20} filtra solo il rumore del Joy1 sul controller; \codepath{TURRET_INPUT_DEADZONE=80} mantiene la sensibilita' storica del Joy2. Queste soglie non sostituiscono \codepath{TRACK_COMMAND_DEADZONE=100} sul tank, applicata anche al risultato del mixing differenziale.
\item \codepath{DebouncedButton} separa livello candidato e livello stabile: un cambiamento vale solo dopo 40 ms continui. Dopo boot/reconnect, \codepath{JoystickReader_requireNeutralBeforeCommands()} richiede 400 ms di quattro assi neutrali e due pulsanti rilasciati.
\item \codepath{udpReady} dice che il socket locale e' stato aperto. \codepath{configureUdpConnection()} controlla \codepath{udp.begin()}, richiede di nuovo neutralita' e marca il primo campione della riconnessione da scartare.
\item \codepath{UdpSender_send()} controlla lunghezza \codepath{snprintf}, \codepath{beginPacket}, byte di \codepath{print} e \codepath{endPacket}. Un errore rimette \codepath{udpReady} a falso e il modulo ritenta la configurazione del socket.
\end{itemize}

\section{Lettura completa dei file: baseline e sorgente corrente}
\subsection{Configurazione PlatformIO del tank}
\textbf{File:} \codepath{tank/platformio.ini}\\
\textbf{Collegamenti audit:} \codepath{A-10, A-14}

\subsubsection{Panoramica}
Questa configurazione seleziona il microcontrollore RA4M1 dell'Arduino Uno R4 WiFi, il framework Arduino e il baud rate del monitor seriale. Non dichiara dipendenze esterne: \codepath{Wire} e \codepath{WiFiS3} arrivano dal framework, mentre \codepath{motorController} e \codepath{PWMController} sono librerie locali sotto \codepath{tank/lib}. Il valore \codepath{monitor_speed=9600} deve restare coerente con \codepath{Serial.begin(9600)} nel firmware.

\subsubsection{Blocchi logici}
\begin{longtable}{p{0.14\textwidth}p{0.24\textwidth}p{0.50\textwidth}}
\toprule
\textbf{Righe} & \textbf{Blocco} & \textbf{Cosa studiare} \\
\midrule
\endhead
1 & Ambiente & Nomina l'ambiente di build usato dal comando PlatformIO. \\
2--5 & Piattaforma e seriale & Sceglie package, board, framework Arduino e velocita' del monitor; una versione platform non fissata puo' cambiare nel tempo. \\
\bottomrule
\end{longtable}

\subsubsection{Funzioni e failure mode}
Non contiene funzioni C++. E' una dichiarazione di build: PlatformIO la legge prima della compilazione per risolvere compilatore, core Arduino, bootloader e librerie.

\subsubsection{Listing completo numerato}
\lstinputlisting[language=]{../tank/platformio.ini}

\subsubsection{Lettura riga per riga}
\begin{longtable}{p{0.10\textwidth}p{0.78\textwidth}}
\toprule
\textbf{Riga} & \textbf{Lettura riga per riga} \\
\midrule
\endhead
1 & Riga sintatticamente significativa del blocco corrente; il suo ruolo e' dettagliato dal contesto e dal listing adiacente. \\
2 & Aggiorna o inizializza uno stato; seguire la variabile permette di vedere l'effetto sul ciclo successivo. \\
3 & Aggiorna o inizializza uno stato; seguire la variabile permette di vedere l'effetto sul ciclo successivo. \\
4 & Aggiorna o inizializza uno stato; seguire la variabile permette di vedere l'effetto sul ciclo successivo. \\
5 & Aggiorna o inizializza uno stato; seguire la variabile permette di vedere l'effetto sul ciclo successivo. \\
\bottomrule
\end{longtable}


\clearpage
\subsection{Coordinatore principale del tank}
\textbf{File:} \codepath{tank/src/main.cpp}\\
\textbf{Collegamenti audit:} \codepath{A-01, A-02, A-03, A-05, A-07, A-10, A-13}

\subsubsection{Panoramica}
\codepath{main.cpp} non implementa tutti i movimenti direttamente: crea lo stato globale, esegue il boot e collega il comando ricevuto ai moduli di cingoli, torretta, elevazione e relay. Il modello Arduino chiama \codepath{setup()} una volta e poi richiama \codepath{loop()} senza fine. Ogni azione e' cooperativa: se una chiamata blocca il processore, anche timeout, relay e movimento attendono il ritorno nel loop.

\subsubsection{Blocchi logici}
\begin{longtable}{p{0.14\textwidth}p{0.24\textwidth}p{0.50\textwidth}}
\toprule
\textbf{Righe} & \textbf{Blocco} & \textbf{Cosa studiare} \\
\midrule
\endhead
1--24 & Dipendenze e contratto & Importa API Arduino/I2C e interfacce locali; il commento documenta il routing dal pacchetto agli attuatori. \\
26--56 & Configurazione e stato persistente & Definisce polarita' del relay, tempi, oggetto PCA9685 e timestamp/stati del relay. Gli oggetti statici esistono per tutta la vita del firmware. \\
58--80 & Ricerca PCA9685 & Sonda tre indirizzi I2C e salva il primo che risponde. Un ACK prova che qualcosa risponde, non che sia necessariamente una PCA9685. \\
82--112 & Relay Fire & Converte lo stato logico in livello elettrico e genera un impulso sul fronte della pressione, con cooldown. \\
114--122 & Stop centrale & Arresta cingoli tramite shield e stepper tramite D2--D5; non spegne i servo, che mantengono l'ultimo PWM. \\
124--149 & Boot & Avvia seriale, relay, stepper, I2C, shield e AP UDP. Il ritardo iniziale e l'assenza di un enable hardware sono importanti per A-01 e A-02. \\
151--182 & Ciclo di controllo & Legge il comando, aggiorna attuatori, applica il fail-safe software e gestisce il fronte Zero. \\
184--208 & Telemetria & Stampa lo stato una volta al secondo; e' utile al debug ma puo' aumentare il jitter del loop. \\
\bottomrule
\end{longtable}

\subsubsection{Funzioni e failure mode}
\textbf{Funzioni.} \codepath{i2cDevicePresent} invia una transazione vuota e interpreta il codice zero come ACK. \codepath{configureShieldAddress} cerca e imposta un indirizzo. \codepath{setFireRelay} applica la polarita' configurata al pin D7. \codepath{updateFireRelay} conserva il fronte precedente, misura impulso e cooldown tramite \codepath{millis}. \codepath{stopAllMotion} delega i due stop ai moduli. \codepath{setup} inizializza hardware e rete; \codepath{loop} e' lo scheduler cooperativo che collega tutte le parti.

\subsubsection{Listing completo numerato}
\lstinputlisting[language=C++]{../tank/src/main.cpp}

\subsubsection{Lettura riga per riga}
\begin{longtable}{p{0.10\textwidth}p{0.78\textwidth}}
\toprule
\textbf{Riga} & \textbf{Lettura riga per riga} \\
\midrule
\endhead
1 & Importa il core Arduino: tipi, GPIO, tempo, map e constrain. \\
2 & Importa l'API I2C Wire usata per parlare con la PCA9685. \\
4 & Include il modulo locale motorController.h e rende visibili la sua interfaccia e i suoi tipi. \\
5 & Include il modulo locale servoTorreta.h e rende visibili la sua interfaccia e i suoi tipi. \\
6 & Include il modulo locale trackController.h e rende visibili la sua interfaccia e i suoi tipi. \\
7 & Include il modulo locale udpReceiver.h e rende visibili la sua interfaccia e i suoi tipi. \\
29 & Configura il pin hardware tramite la macro FIRE\_\allowbreak{}RELAY\_\allowbreak{}PIN; il valore associato deve coincidere con il cablaggio. \\
32 & Definisce la costante di compilazione FIRE\_\allowbreak{}RELAY\_\allowbreak{}ACTIVE\_\allowbreak{}HIGH con valore true. \\
35 & Definisce una durata nominale in millisecondi per FIRE\_\allowbreak{}RELAY\_\allowbreak{}PULSE\_\allowbreak{}MS; viene rispettata solo quando il loop continua a girare. \\
38 & Definisce una durata nominale in millisecondi per FIRE\_\allowbreak{}RELAY\_\allowbreak{}COOLDOWN\_\allowbreak{}MS; viene rispettata solo quando il loop continua a girare. \\
43 & Definisce una durata nominale in millisecondi per COORDINATE\_\allowbreak{}PRINT\_\allowbreak{}INTERVAL\_\allowbreak{}MS; viene rispettata solo quando il loop continua a girare. \\
47 & Dichiara l'API shield senza implementarla: la definizione e' in un file .cpp o piu' avanti. \\
50 & Dichiara stato statico: dura per l'intera esecuzione e, se locale a funzione, conserva valore fra chiamate. \\
53 & Dichiara stato statico: dura per l'intera esecuzione e, se locale a funzione, conserva valore fra chiamate. \\
54 & Dichiara stato statico: dura per l'intera esecuzione e, se locale a funzione, conserva valore fra chiamate. \\
55 & Dichiara stato statico: dura per l'intera esecuzione e, se locale a funzione, conserva valore fra chiamate. \\
56 & Dichiara stato statico: dura per l'intera esecuzione e, se locale a funzione, conserva valore fra chiamate. \\
59 & Definisce la funzione i2cDevicePresent; i parametri costituiscono il suo input e il corpo seguente ne realizza il contratto. \\
60 & Esegue un'operazione I2C sul bus condiviso; il suo esito e la sua durata influenzano il loop cooperativo. \\
61 & Termina la funzione e restituisce al chiamante il valore o il controllo indicato. \\
66 & Definisce la funzione configureShieldAddress; i parametri costituiscono il suo input e il corpo seguente ne realizza il contratto. \\
67 & Dichiara stato statico: dura per l'intera esecuzione e, se locale a funzione, conserva valore fra chiamate. \\
69 & Avvia un ciclo con inizializzazione, condizione di prosecuzione e avanzamento espliciti. \\
70 & Valuta una guardia condizionale prima di eseguire il blocco successivo. \\
71 & Esegue l'istruzione o la chiamata indicata nel contesto della funzione corrente. \\
72 & Emette diagnostica seriale; e' osservabilita' utile ma puo' aggiungere tempo al giro di loop. \\
73 & Emette diagnostica seriale; e' osservabilita' utile ma puo' aggiungere tempo al giro di loop. \\
74 & Termina la funzione e restituisce al chiamante il valore o il controllo indicato. \\
78 & Emette diagnostica seriale; e' osservabilita' utile ma puo' aggiungere tempo al giro di loop. \\
79 & Termina la funzione e restituisce al chiamante il valore o il controllo indicato. \\
83 & Definisce la funzione setFireRelay; i parametri costituiscono il suo input e il corpo seguente ne realizza il contratto. \\
84 & Dichiara o inizializza una variabile: il tipo determina intervallo, durata e semantica del dato. \\
85 & Esegue l'istruzione o la chiamata indicata nel contesto della funzione corrente. \\
87 & Scrive un livello logico sul GPIO e quindi ha un effetto elettrico diretto sul driver/relay. \\
88 & Aggiorna o inizializza uno stato; seguire la variabile permette di vedere l'effetto sul ciclo successivo. \\
93 & Definisce la funzione updateFireRelay; i parametri costituiscono il suo input e il corpo seguente ne realizza il contratto. \\
94 & Dichiara stato statico: dura per l'intera esecuzione e, se locale a funzione, conserva valore fra chiamate. \\
96 & Dichiara o inizializza una variabile: il tipo determina intervallo, durata e semantica del dato. \\
97 & Dichiara o inizializza una variabile: il tipo determina intervallo, durata e semantica del dato. \\
98 & Legge o confronta il contatore temporale Arduino; la sottrazione unsigned e' rollover-safe ma richiede un loop che torni a eseguire. \\
100 & Valuta una guardia condizionale prima di eseguire il blocco successivo. \\
101 & Esegue l'istruzione o la chiamata indicata nel contesto della funzione corrente. \\
102 & Legge o confronta il contatore temporale Arduino; la sottrazione unsigned e' rollover-safe ma richiede un loop che torni a eseguire. \\
103 & Aggiorna o inizializza uno stato; seguire la variabile permette di vedere l'effetto sul ciclo successivo. \\
104 & Aggiorna o inizializza uno stato; seguire la variabile permette di vedere l'effetto sul ciclo successivo. \\
107 & Valuta una guardia condizionale prima di eseguire il blocco successivo. \\
108 & Esegue l'istruzione o la chiamata indicata nel contesto della funzione corrente. \\
111 & Aggiorna o inizializza uno stato; seguire la variabile permette di vedere l'effetto sul ciclo successivo. \\
116 & Definisce la funzione stopAllMotion; i parametri costituiscono il suo input e il corpo seguente ne realizza il contratto. \\
117 & Valuta una guardia condizionale prima di eseguire il blocco successivo. \\
118 & Esegue l'istruzione o la chiamata indicata nel contesto della funzione corrente. \\
121 & Esegue l'istruzione o la chiamata indicata nel contesto della funzione corrente. \\
124 & Definisce la funzione setup; i parametri costituiscono il suo input e il corpo seguente ne realizza il contratto. \\
125 & Emette diagnostica seriale; e' osservabilita' utile ma puo' aggiungere tempo al giro di loop. \\
126 & Blocca intenzionalmente il loop per il tempo indicato; non e' un timer concorrente. \\
129 & Configura elettricamente la direzione del GPIO prima di leggerlo o pilotarlo. \\
130 & Esegue l'istruzione o la chiamata indicata nel contesto della funzione corrente. \\
133 & Esegue l'istruzione o la chiamata indicata nel contesto della funzione corrente. \\
136 & Esegue un'operazione I2C sul bus condiviso; il suo esito e la sua durata influenzano il loop cooperativo. \\
137 & Aggiorna o inizializza uno stato; seguire la variabile permette di vedere l'effetto sul ciclo successivo. \\
138 & Valuta una guardia condizionale prima di eseguire il blocco successivo. \\
140 & Esegue l'istruzione o la chiamata indicata nel contesto della funzione corrente. \\
141 & Esegue l'istruzione o la chiamata indicata nel contesto della funzione corrente. \\
142 & Esegue l'istruzione o la chiamata indicata nel contesto della funzione corrente. \\
143 & Esegue l'istruzione o la chiamata indicata nel contesto della funzione corrente. \\
147 & Esegue l'istruzione o la chiamata indicata nel contesto della funzione corrente. \\
148 & Emette diagnostica seriale; e' osservabilita' utile ma puo' aggiungere tempo al giro di loop. \\
151 & Definisce la funzione loop; i parametri costituiscono il suo input e il corpo seguente ne realizza il contratto. \\
154 & Dichiara o inizializza una variabile: il tipo determina intervallo, durata e semantica del dato. \\
157 & Esegue l'istruzione o la chiamata indicata nel contesto della funzione corrente. \\
160 & Valuta una guardia condizionale prima di eseguire il blocco successivo. \\
161 & Esegue l'istruzione o la chiamata indicata nel contesto della funzione corrente. \\
162 & Esegue l'istruzione o la chiamata indicata nel contesto della funzione corrente. \\
165 & Esegue l'istruzione o la chiamata indicata nel contesto della funzione corrente. \\
169 & Valuta una guardia condizionale prima di eseguire il blocco successivo. \\
170 & Esegue l'istruzione o la chiamata indicata nel contesto della funzione corrente. \\
175 & Dichiara stato statico: dura per l'intera esecuzione e, se locale a funzione, conserva valore fra chiamate. \\
176 & Valuta una guardia condizionale prima di eseguire il blocco successivo. \\
177 & Esegue l'istruzione o la chiamata indicata nel contesto della funzione corrente. \\
178 & Valuta una guardia condizionale prima di eseguire il blocco successivo. \\
179 & Esegue l'istruzione o la chiamata indicata nel contesto della funzione corrente. \\
182 & Aggiorna o inizializza uno stato; seguire la variabile permette di vedere l'effetto sul ciclo successivo. \\
185 & Dichiara stato statico: dura per l'intera esecuzione e, se locale a funzione, conserva valore fra chiamate. \\
186 & Valuta una guardia condizionale prima di eseguire il blocco successivo. \\
187 & Legge o confronta il contatore temporale Arduino; la sottrazione unsigned e' rollover-safe ma richiede un loop che torni a eseguire. \\
189 & Emette diagnostica seriale; e' osservabilita' utile ma puo' aggiungere tempo al giro di loop. \\
190 & Emette diagnostica seriale; e' osservabilita' utile ma puo' aggiungere tempo al giro di loop. \\
191 & Emette diagnostica seriale; e' osservabilita' utile ma puo' aggiungere tempo al giro di loop. \\
192 & Emette diagnostica seriale; e' osservabilita' utile ma puo' aggiungere tempo al giro di loop. \\
193 & Emette diagnostica seriale; e' osservabilita' utile ma puo' aggiungere tempo al giro di loop. \\
194 & Emette diagnostica seriale; e' osservabilita' utile ma puo' aggiungere tempo al giro di loop. \\
195 & Emette diagnostica seriale; e' osservabilita' utile ma puo' aggiungere tempo al giro di loop. \\
196 & Emette diagnostica seriale; e' osservabilita' utile ma puo' aggiungere tempo al giro di loop. \\
197 & Emette diagnostica seriale; e' osservabilita' utile ma puo' aggiungere tempo al giro di loop. \\
198 & Emette diagnostica seriale; e' osservabilita' utile ma puo' aggiungere tempo al giro di loop. \\
199 & Emette diagnostica seriale; e' osservabilita' utile ma puo' aggiungere tempo al giro di loop. \\
200 & Emette diagnostica seriale; e' osservabilita' utile ma puo' aggiungere tempo al giro di loop. \\
201 & Emette diagnostica seriale; e' osservabilita' utile ma puo' aggiungere tempo al giro di loop. \\
202 & Emette diagnostica seriale; e' osservabilita' utile ma puo' aggiungere tempo al giro di loop. \\
203 & Emette diagnostica seriale; e' osservabilita' utile ma puo' aggiungere tempo al giro di loop. \\
204 & Emette diagnostica seriale; e' osservabilita' utile ma puo' aggiungere tempo al giro di loop. \\
205 & Emette diagnostica seriale; e' osservabilita' utile ma puo' aggiungere tempo al giro di loop. \\
206 & Emette diagnostica seriale; e' osservabilita' utile ma puo' aggiungere tempo al giro di loop. \\
\bottomrule
\end{longtable}


\clearpage
\subsection{Interfaccia del ricevitore UDP}
\textbf{File:} \codepath{tank/src/udpReceiver.h}\\
\textbf{Collegamenti audit:} \codepath{A-01, A-08, A-09, A-13, A-14}

\subsubsection{Panoramica}
Questo header separa la rete dal resto dell'applicazione. La struct \codepath{TankCommand} e' il contratto interno: quattro assi sono in scala 0--1023, due booleani rappresentano pulsanti e \codepath{connected} segnala che un pacchetto valido e' arrivato recentemente. Non significa che il controller sia certamente presente o autenticato.

\subsubsection{Blocchi logici}
\begin{longtable}{p{0.14\textwidth}p{0.24\textwidth}p{0.50\textwidth}}
\toprule
\textbf{Righe} & \textbf{Blocco} & \textbf{Cosa studiare} \\
\midrule
\endhead
1--4, 39 & Protezione header & Evita di includere la stessa dichiarazione due volte nella medesima unita' di compilazione. \\
14--31 & Struttura TankCommand & Raccoglie il valore richiesto da ogni attuatore e lo stato del watchdog software. \\
33--37 & API pubblica & Espone avvio dell'AP e lettura dell'ultimo comando in forma indipendente da WiFi/UDP. \\
\bottomrule
\end{longtable}

\subsubsection{Funzioni e failure mode}
\codepath{UdpReceiver_begin} prepara AP e socket. \codepath{UdpReceiver_update} restituisce un \codepath{TankCommand} per valore: il chiamante riceve una copia coerente dell'ultimo stato memorizzato.

\subsubsection{Listing completo numerato}
\lstinputlisting[language=C++]{../tank/src/udpReceiver.h}

\subsubsection{Lettura riga per riga}
\begin{longtable}{p{0.10\textwidth}p{0.78\textwidth}}
\toprule
\textbf{Riga} & \textbf{Lettura riga per riga} \\
\midrule
\endhead
1 & Inizia una include guard: il corpo dell'header verra' elaborato solo se la macro non e' gia' definita. \\
2 & Definisce la macro della include guard: impedisce inclusioni multiple dell'header. \\
4 & Importa il core Arduino: tipi, GPIO, tempo, map e constrain. \\
16 & Dichiara la struct TankCommand: un contenitore di campi che viaggiano insieme fra moduli. \\
18 & Dichiara o inizializza una variabile: il tipo determina intervallo, durata e semantica del dato. \\
19 & Dichiara o inizializza una variabile: il tipo determina intervallo, durata e semantica del dato. \\
22 & Dichiara o inizializza una variabile: il tipo determina intervallo, durata e semantica del dato. \\
23 & Dichiara o inizializza una variabile: il tipo determina intervallo, durata e semantica del dato. \\
26 & Dichiara o inizializza una variabile: il tipo determina intervallo, durata e semantica del dato. \\
27 & Dichiara o inizializza una variabile: il tipo determina intervallo, durata e semantica del dato. \\
30 & Dichiara o inizializza una variabile: il tipo determina intervallo, durata e semantica del dato. \\
34 & Dichiara l'API UdpReceiver\_\allowbreak{}begin senza implementarla: la definizione e' in un file .cpp o piu' avanti. \\
37 & Dichiara l'API UdpReceiver\_\allowbreak{}update senza implementarla: la definizione e' in un file .cpp o piu' avanti. \\
39 & Chiude la include guard o il blocco di precompilazione aperto in precedenza. \\
\bottomrule
\end{longtable}


\clearpage
\subsection{AP WiFi e parser del comando UDP}
\textbf{File:} \codepath{tank/src/udpReceiver.cpp}\\
\textbf{Collegamenti audit:} \codepath{A-01, A-04, A-08, A-09, A-10, A-13, A-14}

\subsubsection{Panoramica}
Il tank opera come access point e ascolta UDP sulla porta 4210. Il protocollo V1 e' testuale e compatto. Il modulo conserva \codepath{lastCommand} fra i giri di loop e lo sostituisce solo quando il parsing riesce. Dopo 500 ms senza pacchetto valido restituisce valori centrati e \codepath{connected=false}; questa e' una protezione logica, non un interruttore elettrico autonomo.

\subsubsection{Blocchi logici}
\begin{longtable}{p{0.14\textwidth}p{0.24\textwidth}p{0.50\textwidth}}
\toprule
\textbf{Righe} & \textbf{Blocco} & \textbf{Cosa studiare} \\
\midrule
\endhead
1--18 & Librerie e contratto V1 & Carica WiFiS3 e WiFiUdp e descrive formato, scala e timeout. \\
20--43 & Parametri e stato statico & Definisce credenziali, porta, buffer, comando sicuro iniziale e timestamp. Il buffer e' grande 64 byte ma la lettura ne usa al massimo 63. \\
45--61 & Avvio AP & Crea l'access point, blocca per sempre se fallisce, aspetta un secondo e apre UDP. \\
63--99 & Ricezione e parsing & Legge un datagramma, lo termina, esegue \codepath{sscanf}, limita assi e aggiorna il timestamp solo se sei campi sono letti. \\
101--112 & Timeout e return & Ripristina valori neutrali se l'eta' dell'ultimo pacchetto supera il limite. \\
\bottomrule
\end{longtable}

\subsubsection{Funzioni e failure mode}
\codepath{UdpReceiver_begin} effettua una sola inizializzazione dell'AP. \codepath{UdpReceiver_update} e' chiamata ad ogni giro: \codepath{parsePacket} chiede al sottosistema WiFi se esiste un datagramma, \codepath{read} copia testo nel buffer e il parser aggiorna lo stato. I fallback locali evitano di usare variabili non inizializzate, ma non rendono il parser una grammatica rigorosa.

\subsubsection{Listing completo numerato}
\lstinputlisting[language=C++]{../tank/src/udpReceiver.cpp}

\subsubsection{Lettura riga per riga}
\begin{longtable}{p{0.10\textwidth}p{0.78\textwidth}}
\toprule
\textbf{Riga} & \textbf{Lettura riga per riga} \\
\midrule
\endhead
1 & Include il modulo locale udpReceiver.h e rende visibili la sua interfaccia e i suoi tipi. \\
3 & Importa l'API WiFi della board corrente. \\
4 & Importa il socket UDP fornito dal framework. \\
21 & Dichiara o inizializza una variabile: il tipo determina intervallo, durata e semantica del dato. \\
22 & Dichiara o inizializza una variabile: il tipo determina intervallo, durata e semantica del dato. \\
25 & Fissa la porta di comunicazione UDP\_\allowbreak{}PORT condivisa dai due firmware. \\
28 & Definisce una durata nominale in millisecondi per CONNECTION\_\allowbreak{}TIMEOUT\_\allowbreak{}MS; viene rispettata solo quando il loop continua a girare. \\
31 & Definisce una soglia di neutralita' del joystick per JOYSTICK\_\allowbreak{}CENTER. \\
33 & Dichiara stato statico: dura per l'intera esecuzione e, se locale a funzione, conserva valore fra chiamate. \\
36 & Dichiara stato statico: dura per l'intera esecuzione e, se locale a funzione, conserva valore fra chiamate. \\
39 & Dichiara stato statico: dura per l'intera esecuzione e, se locale a funzione, conserva valore fra chiamate. \\
40 & Esegue l'istruzione o la chiamata indicata nel contesto della funzione corrente. \\
43 & Dichiara stato statico: dura per l'intera esecuzione e, se locale a funzione, conserva valore fra chiamate. \\
45 & Definisce la funzione UdpReceiver\_\allowbreak{}begin; i parametri costituiscono il suo input e il corpo seguente ne realizza il contratto. \\
46 & Emette diagnostica seriale; e' osservabilita' utile ma puo' aggiungere tempo al giro di loop. \\
49 & Valuta una guardia condizionale prima di eseguire il blocco successivo. \\
50 & Emette diagnostica seriale; e' osservabilita' utile ma puo' aggiungere tempo al giro di loop. \\
51 & Ripete il blocco finche' la condizione resta vera; un loop infinito richiede particolare attenzione al fail-safe. \\
55 & Blocca intenzionalmente il loop per il tempo indicato; non e' un timer concorrente. \\
56 & Invoca il socket UDP; un successo locale non equivale a conferma di ricezione remota. \\
59 & Emette diagnostica seriale; e' osservabilita' utile ma puo' aggiungere tempo al giro di loop. \\
60 & Emette diagnostica seriale; e' osservabilita' utile ma puo' aggiungere tempo al giro di loop. \\
63 & Definisce la funzione UdpReceiver\_\allowbreak{}update; i parametri costituiscono il suo input e il corpo seguente ne realizza il contratto. \\
65 & Dichiara o inizializza una variabile: il tipo determina intervallo, durata e semantica del dato. \\
67 & Valuta una guardia condizionale prima di eseguire il blocco successivo. \\
69 & Dichiara o inizializza una variabile: il tipo determina intervallo, durata e semantica del dato. \\
70 & Valuta una guardia condizionale prima di eseguire il blocco successivo. \\
71 & Aggiorna o inizializza uno stato; seguire la variabile permette di vedere l'effetto sul ciclo successivo. \\
74 & Dichiara o inizializza una variabile: il tipo determina intervallo, durata e semantica del dato. \\
75 & Dichiara o inizializza una variabile: il tipo determina intervallo, durata e semantica del dato. \\
76 & Dichiara o inizializza una variabile: il tipo determina intervallo, durata e semantica del dato. \\
77 & Dichiara o inizializza una variabile: il tipo determina intervallo, durata e semantica del dato. \\
78 & Dichiara o inizializza una variabile: il tipo determina intervallo, durata e semantica del dato. \\
79 & Dichiara o inizializza una variabile: il tipo determina intervallo, durata e semantica del dato. \\
82 & Interpreta il testo ricevuto secondo un formato; il numero di campi letti decide se il frame viene accettato. \\
83 & Esegue l'istruzione o la chiamata indicata nel contesto della funzione corrente. \\
85 & Valuta una guardia condizionale prima di eseguire il blocco successivo. \\
87 & Limita un valore all'intervallo ammesso prima dell'uso successivo. \\
88 & Limita un valore all'intervallo ammesso prima dell'uso successivo. \\
89 & Limita un valore all'intervallo ammesso prima dell'uso successivo. \\
90 & Limita un valore all'intervallo ammesso prima dell'uso successivo. \\
91 & Aggiorna o inizializza uno stato; seguire la variabile permette di vedere l'effetto sul ciclo successivo. \\
92 & Aggiorna o inizializza uno stato; seguire la variabile permette di vedere l'effetto sul ciclo successivo. \\
93 & Aggiorna o inizializza uno stato; seguire la variabile permette di vedere l'effetto sul ciclo successivo. \\
96 & Legge o confronta il contatore temporale Arduino; la sottrazione unsigned e' rollover-safe ma richiede un loop che torni a eseguire. \\
102 & Valuta una guardia condizionale prima di eseguire il blocco successivo. \\
103 & Aggiorna o inizializza uno stato; seguire la variabile permette di vedere l'effetto sul ciclo successivo. \\
104 & Aggiorna o inizializza uno stato; seguire la variabile permette di vedere l'effetto sul ciclo successivo. \\
105 & Aggiorna o inizializza uno stato; seguire la variabile permette di vedere l'effetto sul ciclo successivo. \\
106 & Aggiorna o inizializza uno stato; seguire la variabile permette di vedere l'effetto sul ciclo successivo. \\
107 & Aggiorna o inizializza uno stato; seguire la variabile permette di vedere l'effetto sul ciclo successivo. \\
108 & Aggiorna o inizializza uno stato; seguire la variabile permette di vedere l'effetto sul ciclo successivo. \\
109 & Aggiorna o inizializza uno stato; seguire la variabile permette di vedere l'effetto sul ciclo successivo. \\
112 & Termina la funzione e restituisce al chiamante il valore o il controllo indicato. \\
\bottomrule
\end{longtable}


\clearpage
\subsection{Interfaccia e configurazione dei cingoli}
\textbf{File:} \codepath{tank/src/trackController.h}\\
\textbf{Collegamenti audit:} \codepath{A-11}

\subsubsection{Panoramica}
Il modulo implementa guida differenziale per due motori DC: Joy1 Y decide avanti/indietro e Joy1 X la sterzata. Il lato sinistro e' M3, il destro M1. \codepath{TRACK_COMMAND_DEADZONE=100} e' una protezione finale del tank: viene applicata sia ai due assi prima del mixing, sia ai due comandi dei cingoli dopo \codepath{forward +/- turn}. Non e' la stessa soglia di rumore dell'ESP32. Le inversioni degli assi di guida sono definite soltanto nel controller; qui rimangono soltanto le inversioni dovute al cablaggio dei due motori.

\subsubsection{Blocchi logici}
\begin{longtable}{p{0.14\textwidth}p{0.24\textwidth}p{0.50\textwidth}}
\toprule
\textbf{Righe} & \textbf{Blocco} & \textbf{Cosa studiare} \\
\midrule
\endhead
1--24 & Dipendenze e modello differenziale & Descrive assi, motori M3/M1 e formule concettuali sinistra=forward+turn, destra=forward-turn. \\
26--36 & Scala e sicurezza comando & Fissa centro 512 e \codepath{TRACK_COMMAND_DEADZONE=100}, distinta dalla deadzone fisica del controller. \\
38--60 & Porte, PWM e trim & Seleziona le porte motore e imposta soglie minimo/massimo e compensazioni per i due cingoli. \\
62--68 & Sterzo lineare e inversioni hardware & Stabilisce che Joy1 X entra direttamente nel mixing, senza modellazione aggiuntiva o soglie intermedie, e definisce i soli flag di inversione dei motori. \\
70--79 & API & Dichiara inizializzazione, aggiornamento e stop. \\
\bottomrule
\end{longtable}

\subsubsection{Funzioni e failure mode}
Le tre API pubbliche ricevono o conservano un riferimento al \codepath{MotorController}; l'header non compie I/O ma definisce i limiti che governano la successiva scrittura I2C. \codepath{TrackController_stop()} usa \codepath{DCbrake()} per portare il PWM a zero: il suo nome non dimostra un taglio fisico dell'alimentazione o una frenatura meccanica misurata.

\subsubsection{Listing completo numerato}
\lstinputlisting[language=C++]{../tank/src/trackController.h}

\subsubsection{Lettura riga per riga}
\begin{longtable}{p{0.10\textwidth}p{0.78\textwidth}}
\toprule
\textbf{Riga} & \textbf{Lettura riga per riga} \\
\midrule
\endhead
1 & Inizia una include guard: il corpo dell'header verra' elaborato solo se la macro non e' gia' definita. \\
2 & Definisce la macro della include guard: impedisce inclusioni multiple dell'header. \\
4 & Importa il core Arduino: tipi, GPIO, tempo, map e constrain. \\
6 & Include il modulo locale motorController.h e rende visibili la sua interfaccia e i suoi tipi. \\
28--30 & Documenta e definisce \codepath{DRIVE_JOYSTICK_CENTER=512}: e' il centro della scala UDP, non una misura diretta dell'ADC a 12 bit. \\
32--36 & Definisce \codepath{TRACK_COMMAND_DEADZONE=100}. Il commento chiarisce i due punti di applicazione e perche' non va resa uguale a \codepath{DRIVE_INPUT_DEADZONE} dell'ESP32. \\
40--41 & Associa M3 al cingolo sinistro e M1 al destro; un errore qui scambia il lato azionato. \\
51--54 & Imposta PWM minimo e massimo per ciascun lato. \codepath{MAX_SPEED} arriva dal driver; la partenza reale va verificata con il carico meccanico. \\
59--60 & Imposta il trim percentuale finale: 100 non corregge, valori diversi compensano differenze tra i motori ma possono ridurre il PWM minimo effettivo. \\
62--63 & Documenta lo sterzo lineare: \codepath{driveX} entra direttamente nelle formule del mixing, senza percentuali o soglie intermedie da tarare. \\
65--68 & Riserva le inversioni a M3/M1 per il cablaggio dei motori. Non esistono inversioni \codepath{DRIVE_X/Y} nel tank: Joy1 e' orientato una sola volta nell'ESP32. \\
71 & Dichiara \codepath{TrackController_begin()}, che memorizza il riferimento al driver e parte fermo. \\
74 & Dichiara \codepath{TrackController_update()}, chiamata dal loop del tank con \codepath{driveX} e \codepath{driveY} del frame autorizzato. \\
77 & Dichiara \codepath{TrackController_stop()}, che richiede \codepath{DCbrake()} su entrambi i cingoli. \\
79 & Chiude la include guard. \\
\bottomrule
\end{longtable}


\clearpage
\subsection{Mixing differenziale e PWM dei cingoli}
\textbf{File:} \codepath{tank/src/trackController.cpp}\\
\textbf{Collegamenti audit:} \codepath{A-11}

\subsubsection{Panoramica}
Il file trasforma il protocollo 0--1023 in due comandi con segno. \codepath{decodeDriveAxis()} centra e filtra ciascun asse sul tank; lo sterzo entra linearmente nel mixing \codepath{forward +/- turn}; \codepath{applyTrackMotorCommand()} filtra di nuovo il comando di ogni singolo cingolo e lo traduce in direzione e PWM. Gli input sono limitati e il prodotto PWM per trim usa \codepath{long}, per cui non emerge un overflow pratico nei range correnti. L'orientamento di Joy1 non viene modificato qui: e' gia' normalizzato dall'ESP32.

\subsubsection{Blocchi logici}
\begin{longtable}{p{0.14\textwidth}p{0.24\textwidth}p{0.50\textwidth}}
\toprule
\textbf{Righe} & \textbf{Blocco} & \textbf{Cosa studiare} \\
\midrule
\endhead
14--24 & Driver, dead-time e vincoli & Conserva il puntatore al driver, fissa i 30 ms e verifica a compilazione che la deadzone resti interna alla scala joystick. \\
26--51 & Stato per cingolo & Memorizza ultimo verso richiesto, verso prima dello stop, comando pendente e tempo; sceglie lo stato di M3 o M1. \\
53--89 & Decodifica e stop & \codepath{decodeDriveAxis()} applica il filtro pre-mixing; \codepath{stopTrackMotor()} frena e conserva il contesto necessario al dead-time. \\
91--97 & PWM lineare & Mappa direttamente la magnitudine del comando nel PWM minimo/massimo configurato. \\
99--182 & Applicazione a un motore & Esegue filtro post-mixing, PWM, inversione di cablaggio e macchina di stato per il cambio verso. \\
184--213 & API pubbliche & Associa il driver, calcola il mixing lineare, satura i risultati e arresta entrambi i lati. \\
\bottomrule
\end{longtable}

\subsubsection{Funzioni e failure mode}
\codepath{stateForTrackMotor()} seleziona lo stato statico di M3 o M1. \codepath{decodeDriveAxis()} restituisce un comando firmato e applica il primo filtro \codepath{TRACK_COMMAND_DEADZONE}. \codepath{stopTrackMotor()} richiede \codepath{DCbrake()} e, se il motore era in marcia, conserva verso e \codepath{millis()} dello stop. \codepath{mapLinearTrackSpeed()} calcola il PWM in modo lineare. \codepath{applyTrackMotorCommand()} e' la macchina di stato: il filtro post-mixing e' necessario perche' \codepath{forward +/- turn} puo' lasciare piccolo un solo lato; durante i 30 ms conserva solo l'ultimo comando. Se il joystick passa dal centro al verso opposto, usa \codepath{directionBeforeStop} e \codepath{stoppedAt} per non saltare la pausa. \codepath{TrackController_update()} esegue direttamente il mixing lineare, senza modellare ulteriormente \codepath{turn}.

\subsubsection{Listing completo numerato}
\lstinputlisting[language=C++]{../tank/src/trackController.cpp}

\subsubsection{Lettura riga per riga}
\begin{longtable}{p{0.10\textwidth}p{0.78\textwidth}}
\toprule
\textbf{Riga} & \textbf{Lettura riga per riga} \\
\midrule
\endhead
1 & Include il modulo locale trackController.h e rende visibili la sua interfaccia e i suoi tipi. \\
14--15 & Dichiara il puntatore statico \codepath{motorController}. Resta valido per l'intera esecuzione ma vale \codepath{nullptr} finche' \codepath{TrackController_begin()} non riceve il driver dal main. \\
17--20 & Definisce \codepath{TRACK_REVERSE_DEAD_TIME_MS=30}. E' un intervallo software misurato con \codepath{millis()}, non una garanzia che il ponte H sia fisicamente diseccitato. \\
22--24 & \codepath{static_assert} ferma la compilazione se \codepath{TRACK_COMMAND_DEADZONE} non resta strettamente dentro la semiscala joystick, lasciando almeno un comando positivo e uno negativo. \\
27--37 & Dichiara \codepath{TrackMotorState}. \codepath{lastRequestedDirection} e' l'ultimo verso passato a \codepath{DCrun()}, non un feedback fisico; \codepath{directionBeforeStop} e \codepath{stoppedAt} sopravvivono allo stop; i campi \codepath{pending*} memorizzano l'ultimo comando durante la pausa. \\
39--40 & Inizializza stati indipendenti per M3 sinistro e M1 destro. Nessuna allocazione dinamica avviene nel loop. \\
42--52 & \codepath{stateForTrackMotor()} converte il numero della porta nella struct corretta e restituisce \codepath{nullptr} per una porta non prevista, evitando di dereferenziare stato inesistente. \\
54--58 & \codepath{decodeDriveAxis()} limita l'input UDP a 0--1023 e sottrae 512, ottenendo circa -512..+511. Non inverte Joy1: l'inversione dell'asse appartiene all'ESP32. \\
60--66 & Applica il primo filtro \codepath{TRACK_COMMAND_DEADZONE=100}; un sender UDP alternativo non puo' quindi generare piccoli comandi continui prima del mixing. \\
69--76 & \codepath{stopTrackMotor()} controlla il driver, trova lo stato e invia \codepath{DCbrake()} al motore. La chiamata porta il PWM a zero secondo l'implementazione del driver. \\
79--89 & Se era attivo un verso, salva verso e istante dello stop; azzera solo il verso correntemente richiesto e l'eventuale comando pendente. Questa conservazione rende sicuro il percorso centro -> verso opposto. \\
93--98 & \codepath{mapLinearTrackSpeed()} limita la magnitudine alla porzione utile e la mappa tra PWM minimo e massimo. Un PWM minimo sufficiente a far partire il motore resta una misura hardware. \\
99--114 & \codepath{applyTrackMotorCommand()} riceve porta, inversione di cablaggio, comando firmato e limiti PWM; controlla puntatore/stato e limita tutti gli input prima dell'I/O. \\
116--121 & Applica il secondo filtro da 100 dopo il mixing. E' distinto ma necessario: \codepath{forward+turn} o \codepath{forward-turn} possono produrre un residuo piccolo su un solo cingolo. \\
123--132 & Converte magnitudine in PWM lineare, usa \codepath{long} nel prodotto con il trim, ricava \codepath{FORWARD/BACKWARD} e applica solo l'inversione elettrica del singolo motore. \\
134--150 & Durante un'inversione gia' in attesa aggiorna solo l'ultimo comando pendente. Finche' \codepath{now-stoppedAt < 30}, non richiama \codepath{DCrun()}; allo scadere applica il comando piu' recente. La sottrazione unsigned resta valida al rollover di \codepath{millis()}. \\
153--163 & Se il motore e' in marcia e il verso cambia, invia \codepath{DCbrake()}, salva il verso precedente e avvia la pausa. Non inverte direttamente il ponte H. \\
165--174 & Copre il caso prima assente: dopo joystick al centro, una richiesta immediata del verso opposto confronta \codepath{directionBeforeStop} e aspetta i 30 ms rimanenti invece di ripartire subito. \\
176--182 & Riscrive \codepath{DCrun()} anche a comando invariato. Il commento vieta un caching ingenuo: una transazione I2C periodica aiuta il controllo di salute a osservare un guasto del bus. \\
184--190 & \codepath{TrackController_begin()} memorizza il riferimento al driver gia' inizializzato dal main e avvia il modulo con i cingoli fermati. \\
192--208 & \codepath{TrackController_update()} decodifica Y e X, calcola direttamente \codepath{left=forward+turn} e \codepath{right=forward-turn}, limita i risultati e li invia alle due macchine di stato. \\
210--212 & \codepath{TrackController_stop()} richiede lo stop di entrambi i lati; non disalimenta materialmente la shield. \\
\bottomrule
\end{longtable}


\clearpage
\subsection{Interfaccia torretta ed elevazione}
\textbf{File:} \codepath{tank/src/servoTorreta.h}\\
\textbf{Collegamenti audit:} \codepath{A-07, A-11}

\subsubsection{Panoramica}
Il file riunisce due attuatori diversi: la rotazione orizzontale usa quattro GPIO diretti verso un driver stepper, mentre l'elevazione usa due servo tramite PCA9685. Le API distinguono \codepath{StepperTorretta} e \codepath{ServoTorretta}; il nome comune non implica lo stesso hardware.

\subsubsection{Blocchi logici}
\begin{longtable}{p{0.14\textwidth}p{0.24\textwidth}p{0.50\textwidth}}
\toprule
\textbf{Righe} & \textbf{Blocco} & \textbf{Cosa studiare} \\
\midrule
\endhead
1--21 & Due sottosistemi & Introduce stepper D2--D5 e servo S5/S6 come percorsi separati. \\
23--38 & Costanti stepper & Definisce passi per giro, pin, intervallo e deadzone. \\
40--68 & API stepper e servo & Dichiara operazioni di avvio, movimento, zero, stop e lettura dello stato. \\
\bottomrule
\end{longtable}

\subsubsection{Funzioni e failure mode}
Per lo stepper, Zero resetta soltanto il conteggio software. Per i servo, \codepath{ServoTorretta_setZero} delega a una posizione fisica minima. Questa differenza e' fondamentale per capire A-07.

\subsubsection{Listing completo numerato}
\lstinputlisting[language=C++]{../tank/src/servoTorreta.h}

\subsubsection{Lettura riga per riga}
\begin{longtable}{p{0.10\textwidth}p{0.78\textwidth}}
\toprule
\textbf{Riga} & \textbf{Lettura riga per riga} \\
\midrule
\endhead
1 & Inizia una include guard: il corpo dell'header verra' elaborato solo se la macro non e' gia' definita. \\
2 & Definisce la macro della include guard: impedisce inclusioni multiple dell'header. \\
4 & Importa il core Arduino: tipi, GPIO, tempo, map e constrain. \\
6 & Include il modulo locale motorController.h e rende visibili la sua interfaccia e i suoi tipi. \\
26 & Definisce un limite o una conversione di posizione per STEPS\_\allowbreak{}PER\_\allowbreak{}REV. \\
29 & Configura il pin hardware tramite la macro TURRET\_\allowbreak{}DRIVER\_\allowbreak{}A\_\allowbreak{}PIN; il valore associato deve coincidere con il cablaggio. \\
30 & Configura il pin hardware tramite la macro TURRET\_\allowbreak{}DRIVER\_\allowbreak{}B\_\allowbreak{}PIN; il valore associato deve coincidere con il cablaggio. \\
31 & Configura il pin hardware tramite la macro TURRET\_\allowbreak{}DRIVER\_\allowbreak{}C\_\allowbreak{}PIN; il valore associato deve coincidere con il cablaggio. \\
32 & Configura il pin hardware tramite la macro TURRET\_\allowbreak{}DRIVER\_\allowbreak{}D\_\allowbreak{}PIN; il valore associato deve coincidere con il cablaggio. \\
35 & Definisce una durata nominale in millisecondi per TURRET\_\allowbreak{}STEP\_\allowbreak{}INTERVAL\_\allowbreak{}MS; viene rispettata solo quando il loop continua a girare. \\
38 & Definisce una soglia di neutralita' del joystick per TURRET\_\allowbreak{}JOYSTICK\_\allowbreak{}DEADZONE. \\
41 & Dichiara l'API StepperTorretta\_\allowbreak{}begin senza implementarla: la definizione e' in un file .cpp o piu' avanti. \\
44 & Dichiara l'API StepperTorretta\_\allowbreak{}updateJoystick senza implementarla: la definizione e' in un file .cpp o piu' avanti. \\
47 & Dichiara l'API StepperTorretta\_\allowbreak{}setZero senza implementarla: la definizione e' in un file .cpp o piu' avanti. \\
50 & Dichiara l'API StepperTorretta\_\allowbreak{}stop senza implementarla: la definizione e' in un file .cpp o piu' avanti. \\
53 & Dichiara l'API StepperTorretta\_\allowbreak{}getAngle senza implementarla: la definizione e' in un file .cpp o piu' avanti. \\
56 & Dichiara l'API ServoTorretta\_\allowbreak{}begin senza implementarla: la definizione e' in un file .cpp o piu' avanti. \\
59 & Dichiara l'API ServoTorretta\_\allowbreak{}updateJoystick senza implementarla: la definizione e' in un file .cpp o piu' avanti. \\
62 & Dichiara l'API ServoTorretta\_\allowbreak{}setAngle senza implementarla: la definizione e' in un file .cpp o piu' avanti. \\
65 & Dichiara l'API ServoTorretta\_\allowbreak{}setZero senza implementarla: la definizione e' in un file .cpp o piu' avanti. \\
68 & Dichiara l'API ServoTorretta\_\allowbreak{}getAngle senza implementarla: la definizione e' in un file .cpp o piu' avanti. \\
70 & Chiude la include guard o il blocco di precompilazione aperto in precedenza. \\
\bottomrule
\end{longtable}


\clearpage
\subsection{Sequenza stepper e controllo dei servo}
\textbf{File:} \codepath{tank/src/servoTorreta.cpp}\\
\textbf{Collegamenti audit:} \codepath{A-06, A-07, A-11}

\subsubsection{Panoramica}
La rotazione conserva fase e conteggio software, e alimenta una sola bobina alla volta in sequenza A-B-C-D. L'elevazione conserva un angolo e lo aggiorna di un grado ogni 30 ms. Nessuno dei due percorsi dispone di feedback: i limiti e le coordinate sono creduti dal software, non misurati dalla meccanica.

\subsubsection{Blocchi logici}
\begin{longtable}{p{0.14\textwidth}p{0.24\textwidth}p{0.50\textwidth}}
\toprule
\textbf{Righe} & \textbf{Blocco} & \textbf{Cosa studiare} \\
\midrule
\endhead
18--53 & Limiti e stato & Definisce angoli, canali, intervalli, conversione step/gradi e variabili persistenti. \\
55--86 & Scrittura di fase & Attiva esattamente una combinazione A/B/C/D e segna il driver come abilitato. \\
88--110 & Step destro/sinistro & Controlla limite logico, aggiorna indice modulo quattro e cambia conteggio. \\
112--167 & Lifecycle stepper & Configura pin, limita ritmo con \codepath{millis}, resetta soltanto lo stato logico, spegne bobine e calcola angolo. \\
169--216 & Lifecycle servo & Memorizza shield, aggiorna lentamente angolo e invia coppia specchiata ai due canali. \\
\bottomrule
\end{longtable}

\subsubsection{Funzioni e failure mode}
\codepath{writeTurretStep} porta fisicamente i quattro pin al pattern scelto. \codepath{stepRight} e \codepath{stepLeft} cambiano fase e posizione, ma non possono sapere se il motore ha realmente girato. \codepath{StepperTorretta_updateJoystick} fa al massimo uno step per giro del loop se sono trascorsi 2 ms. \codepath{ServoTorretta_setAngle} limita 0--47 e comanda la coppia. A 2 ms il massimo teorico e' 500 step/s, circa 14,6 rpm assumendo 2048 step/giro; WiFi/I2C/seriale possono solo ridurlo.

\subsubsection{Listing completo numerato}
\lstinputlisting[language=C++]{../tank/src/servoTorreta.cpp}

\subsubsection{Lettura riga per riga}
\begin{longtable}{p{0.10\textwidth}p{0.78\textwidth}}
\toprule
\textbf{Riga} & \textbf{Lettura riga per riga} \\
\midrule
\endhead
1 & Include il modulo locale servoTorreta.h e rende visibili la sua interfaccia e i suoi tipi. \\
21 & Definisce un limite o una conversione di posizione per TURRET\_\allowbreak{}MIN\_\allowbreak{}ANGLE. \\
22 & Definisce un limite o una conversione di posizione per TURRET\_\allowbreak{}MAX\_\allowbreak{}ANGLE. \\
26 & Definisce un limite o una conversione di posizione per ELEVATION\_\allowbreak{}MIN\_\allowbreak{}ANGLE. \\
27 & Definisce un limite o una conversione di posizione per ELEVATION\_\allowbreak{}MAX\_\allowbreak{}ANGLE. \\
28 & Definisce la costante di compilazione ELEVATION\_\allowbreak{}MIRROR\_\allowbreak{}BASE con valore 90. \\
29 & Definisce la costante di compilazione ELEVATION\_\allowbreak{}SERVO\_\allowbreak{}A con valore S5. \\
30 & Definisce la costante di compilazione ELEVATION\_\allowbreak{}SERVO\_\allowbreak{}B con valore S6. \\
33 & Definisce una soglia di neutralita' del joystick per JOYSTICK\_\allowbreak{}CENTER. \\
36 & Definisce una durata nominale in millisecondi per SERVO\_\allowbreak{}INTERVAL\_\allowbreak{}MS; viene rispettata solo quando il loop continua a girare. \\
39 & Definisce un limite o una conversione di posizione per TURRET\_\allowbreak{}MIN\_\allowbreak{}STEPS. \\
40 & Definisce un limite o una conversione di posizione per TURRET\_\allowbreak{}MAX\_\allowbreak{}STEPS. \\
43 & Dichiara stato statico: dura per l'intera esecuzione e, se locale a funzione, conserva valore fra chiamate. \\
46 & Dichiara stato statico: dura per l'intera esecuzione e, se locale a funzione, conserva valore fra chiamate. \\
47 & Dichiara stato statico: dura per l'intera esecuzione e, se locale a funzione, conserva valore fra chiamate. \\
48 & Dichiara stato statico: dura per l'intera esecuzione e, se locale a funzione, conserva valore fra chiamate. \\
49 & Dichiara stato statico: dura per l'intera esecuzione e, se locale a funzione, conserva valore fra chiamate. \\
52 & Dichiara stato statico: dura per l'intera esecuzione e, se locale a funzione, conserva valore fra chiamate. \\
53 & Dichiara stato statico: dura per l'intera esecuzione e, se locale a funzione, conserva valore fra chiamate. \\
57 & Definisce la funzione writeTurretStep; i parametri costituiscono il suo input e il corpo seguente ne realizza il contratto. \\
58 & Apre una selezione multi-ramo basata sul valore dell'espressione. \\
59 & Seleziona questo ramo dello switch quando il valore confrontato coincide con il caso indicato. \\
60 & Scrive un livello logico sul GPIO e quindi ha un effetto elettrico diretto sul driver/relay. \\
61 & Scrive un livello logico sul GPIO e quindi ha un effetto elettrico diretto sul driver/relay. \\
62 & Scrive un livello logico sul GPIO e quindi ha un effetto elettrico diretto sul driver/relay. \\
63 & Scrive un livello logico sul GPIO e quindi ha un effetto elettrico diretto sul driver/relay. \\
64 & Esce dal ramo corrente dello switch evitando il passaggio al caso successivo. \\
65 & Seleziona questo ramo dello switch quando il valore confrontato coincide con il caso indicato. \\
66 & Scrive un livello logico sul GPIO e quindi ha un effetto elettrico diretto sul driver/relay. \\
67 & Scrive un livello logico sul GPIO e quindi ha un effetto elettrico diretto sul driver/relay. \\
68 & Scrive un livello logico sul GPIO e quindi ha un effetto elettrico diretto sul driver/relay. \\
69 & Scrive un livello logico sul GPIO e quindi ha un effetto elettrico diretto sul driver/relay. \\
70 & Esce dal ramo corrente dello switch evitando il passaggio al caso successivo. \\
71 & Seleziona questo ramo dello switch quando il valore confrontato coincide con il caso indicato. \\
72 & Scrive un livello logico sul GPIO e quindi ha un effetto elettrico diretto sul driver/relay. \\
73 & Scrive un livello logico sul GPIO e quindi ha un effetto elettrico diretto sul driver/relay. \\
74 & Scrive un livello logico sul GPIO e quindi ha un effetto elettrico diretto sul driver/relay. \\
75 & Scrive un livello logico sul GPIO e quindi ha un effetto elettrico diretto sul driver/relay. \\
76 & Esce dal ramo corrente dello switch evitando il passaggio al caso successivo. \\
77 & Seleziona questo ramo dello switch quando il valore confrontato coincide con il caso indicato. \\
78 & Scrive un livello logico sul GPIO e quindi ha un effetto elettrico diretto sul driver/relay. \\
79 & Scrive un livello logico sul GPIO e quindi ha un effetto elettrico diretto sul driver/relay. \\
80 & Scrive un livello logico sul GPIO e quindi ha un effetto elettrico diretto sul driver/relay. \\
81 & Scrive un livello logico sul GPIO e quindi ha un effetto elettrico diretto sul driver/relay. \\
82 & Esce dal ramo corrente dello switch evitando il passaggio al caso successivo. \\
85 & Aggiorna o inizializza uno stato; seguire la variabile permette di vedere l'effetto sul ciclo successivo. \\
89 & Definisce la funzione stepRight; i parametri costituiscono il suo input e il corpo seguente ne realizza il contratto. \\
90 & Valuta una guardia condizionale prima di eseguire il blocco successivo. \\
91 & Esegue l'istruzione o la chiamata indicata nel contesto della funzione corrente. \\
92 & Termina la funzione e restituisce al chiamante il valore o il controllo indicato. \\
95 & Aggiorna o inizializza uno stato; seguire la variabile permette di vedere l'effetto sul ciclo successivo. \\
96 & Esegue l'istruzione o la chiamata indicata nel contesto della funzione corrente. \\
97 & Esegue l'istruzione o la chiamata indicata nel contesto della funzione corrente. \\
101 & Definisce la funzione stepLeft; i parametri costituiscono il suo input e il corpo seguente ne realizza il contratto. \\
102 & Valuta una guardia condizionale prima di eseguire il blocco successivo. \\
103 & Esegue l'istruzione o la chiamata indicata nel contesto della funzione corrente. \\
104 & Termina la funzione e restituisce al chiamante il valore o il controllo indicato. \\
107 & Aggiorna o inizializza uno stato; seguire la variabile permette di vedere l'effetto sul ciclo successivo. \\
108 & Esegue l'istruzione o la chiamata indicata nel contesto della funzione corrente. \\
109 & Esegue l'istruzione o la chiamata indicata nel contesto della funzione corrente. \\
112 & Definisce la funzione StepperTorretta\_\allowbreak{}begin; i parametri costituiscono il suo input e il corpo seguente ne realizza il contratto. \\
114 & Configura elettricamente la direzione del GPIO prima di leggerlo o pilotarlo. \\
115 & Configura elettricamente la direzione del GPIO prima di leggerlo o pilotarlo. \\
116 & Configura elettricamente la direzione del GPIO prima di leggerlo o pilotarlo. \\
117 & Configura elettricamente la direzione del GPIO prima di leggerlo o pilotarlo. \\
120 & Aggiorna o inizializza uno stato; seguire la variabile permette di vedere l'effetto sul ciclo successivo. \\
121 & Esegue l'istruzione o la chiamata indicata nel contesto della funzione corrente. \\
122 & Aggiorna o inizializza uno stato; seguire la variabile permette di vedere l'effetto sul ciclo successivo. \\
123 & Aggiorna o inizializza uno stato; seguire la variabile permette di vedere l'effetto sul ciclo successivo. \\
126 & Definisce la funzione StepperTorretta\_\allowbreak{}updateJoystick; i parametri costituiscono il suo input e il corpo seguente ne realizza il contratto. \\
127 & Legge o confronta il contatore temporale Arduino; la sottrazione unsigned e' rollover-safe ma richiede un loop che torni a eseguire. \\
130 & Valuta una guardia condizionale prima di eseguire il blocco successivo. \\
131 & Termina la funzione e restituisce al chiamante il valore o il controllo indicato. \\
134 & Aggiorna o inizializza uno stato; seguire la variabile permette di vedere l'effetto sul ciclo successivo. \\
138 & Valuta una guardia condizionale prima di eseguire il blocco successivo. \\
139 & Esegue l'istruzione o la chiamata indicata nel contesto della funzione corrente. \\
140 & Prosegue la catena condizionale: questo ramo e' valutato solo se quelli precedenti non sono stati scelti. \\
141 & Esegue l'istruzione o la chiamata indicata nel contesto della funzione corrente. \\
142 & Apre l'alternativa della condizione precedente. \\
143 & Esegue l'istruzione o la chiamata indicata nel contesto della funzione corrente. \\
147 & Definisce la funzione StepperTorretta\_\allowbreak{}setZero; i parametri costituiscono il suo input e il corpo seguente ne realizza il contratto. \\
149 & Aggiorna o inizializza uno stato; seguire la variabile permette di vedere l'effetto sul ciclo successivo. \\
152 & Definisce la funzione StepperTorretta\_\allowbreak{}stop; i parametri costituiscono il suo input e il corpo seguente ne realizza il contratto. \\
153 & Valuta una guardia condizionale prima di eseguire il blocco successivo. \\
154 & Termina la funzione e restituisce al chiamante il valore o il controllo indicato. \\
158 & Scrive un livello logico sul GPIO e quindi ha un effetto elettrico diretto sul driver/relay. \\
159 & Scrive un livello logico sul GPIO e quindi ha un effetto elettrico diretto sul driver/relay. \\
160 & Scrive un livello logico sul GPIO e quindi ha un effetto elettrico diretto sul driver/relay. \\
161 & Scrive un livello logico sul GPIO e quindi ha un effetto elettrico diretto sul driver/relay. \\
162 & Aggiorna o inizializza uno stato; seguire la variabile permette di vedere l'effetto sul ciclo successivo. \\
165 & Definisce la funzione StepperTorretta\_\allowbreak{}getAngle; i parametri costituiscono il suo input e il corpo seguente ne realizza il contratto. \\
166 & Termina la funzione e restituisce al chiamante il valore o il controllo indicato. \\
169 & Definisce la funzione ServoTorretta\_\allowbreak{}begin; i parametri costituiscono il suo input e il corpo seguente ne realizza il contratto. \\
171 & Aggiorna o inizializza uno stato; seguire la variabile permette di vedere l'effetto sul ciclo successivo. \\
172 & Esegue l'istruzione o la chiamata indicata nel contesto della funzione corrente. \\
175 & Definisce la funzione ServoTorretta\_\allowbreak{}updateJoystick; i parametri costituiscono il suo input e il corpo seguente ne realizza il contratto. \\
176 & Legge o confronta il contatore temporale Arduino; la sottrazione unsigned e' rollover-safe ma richiede un loop che torni a eseguire. \\
179 & Valuta una guardia condizionale prima di eseguire il blocco successivo. \\
180 & Termina la funzione e restituisce al chiamante il valore o il controllo indicato. \\
183 & Aggiorna o inizializza uno stato; seguire la variabile permette di vedere l'effetto sul ciclo successivo. \\
185 & Dichiara o inizializza una variabile: il tipo determina intervallo, durata e semantica del dato. \\
188 & Valuta una guardia condizionale prima di eseguire il blocco successivo. \\
189 & Esegue l'istruzione o la chiamata indicata nel contesto della funzione corrente. \\
190 & Prosegue la catena condizionale: questo ramo e' valutato solo se quelli precedenti non sono stati scelti. \\
191 & Esegue l'istruzione o la chiamata indicata nel contesto della funzione corrente. \\
194 & Valuta una guardia condizionale prima di eseguire il blocco successivo. \\
195 & Esegue l'istruzione o la chiamata indicata nel contesto della funzione corrente. \\
199 & Definisce la funzione ServoTorretta\_\allowbreak{}setAngle; i parametri costituiscono il suo input e il corpo seguente ne realizza il contratto. \\
200 & Valuta una guardia condizionale prima di eseguire il blocco successivo. \\
201 & Termina la funzione e restituisce al chiamante il valore o il controllo indicato. \\
205 & Limita un valore all'intervallo ammesso prima dell'uso successivo. \\
206 & Delega il comando al controller shield tramite il puntatore precedentemente inizializzato. \\
207 & Esegue l'istruzione o la chiamata indicata nel contesto della funzione corrente. \\
210 & Definisce la funzione ServoTorretta\_\allowbreak{}setZero; i parametri costituiscono il suo input e il corpo seguente ne realizza il contratto. \\
211 & Esegue l'istruzione o la chiamata indicata nel contesto della funzione corrente. \\
214 & Definisce la funzione ServoTorretta\_\allowbreak{}getAngle; i parametri costituiscono il suo input e il corpo seguente ne realizza il contratto. \\
215 & Termina la funzione e restituisce al chiamante il valore o il controllo indicato. \\
\bottomrule
\end{longtable}


\clearpage
\subsection{Interfaccia semantica della shield motori}
\textbf{File:} \codepath{tank/lib/motorController/motorController.h}\\
\textbf{Collegamenti audit:} \codepath{A-02, A-11}

\subsubsection{Panoramica}
Questa classe nasconde i registri PCA9685 dietro operazioni leggibili come \codepath{DCrun} e \codepath{servoTurn}. Le macro mappano porte M1--M4 e S1--S8 a canali PCA. \codepath{MAX_SPEED=4096} rappresenta anche il bit full-on/full-off del protocollo PCA, non un normale valore PWM fra 0 e 4095.

\subsubsection{Blocchi logici}
\begin{longtable}{p{0.14\textwidth}p{0.24\textwidth}p{0.50\textwidth}}
\toprule
\textbf{Righe} & \textbf{Blocco} & \textbf{Cosa studiare} \\
\midrule
\endhead
1--48 & Identificatori e mappa hardware & Definisce motori, canali di direzione/velocita', direzioni, range PWM e connettori servo. \\
50--84 & Classe e API & Dichiara costruttore, operazioni DC/servo e helper privati verso il driver basso livello. \\
\bottomrule
\end{longtable}

\subsubsection{Funzioni e failure mode}
Il costruttore prende frequenza e indirizzo. Le API \codepath{DCrun} e \codepath{DCbrake} esprimono il controllo motore. Le API \codepath{servoTurn} e \codepath{servoPairTurn} esprimono un impulso servo. Gli helper privati decidono se usare PWM normale o flag full-on/full-off.

\subsubsection{Listing completo numerato}
\lstinputlisting[language=C++]{../tank/lib/motorController/motorController.h}

\subsubsection{Lettura riga per riga}
\begin{longtable}{p{0.10\textwidth}p{0.78\textwidth}}
\toprule
\textbf{Riga} & \textbf{Lettura riga per riga} \\
\midrule
\endhead
1 & Inizia una include guard: il corpo dell'header verra' elaborato solo se la macro non e' gia' definita. \\
2 & Definisce la macro della include guard: impedisce inclusioni multiple dell'header. \\
4 & Include il modulo locale PWMController.h e rende visibili la sua interfaccia e i suoi tipi. \\
18 & Definisce la costante di compilazione M1 con valore 1. \\
19 & Definisce la costante di compilazione M2 con valore 2. \\
20 & Definisce la costante di compilazione M3 con valore 3. \\
21 & Definisce la costante di compilazione M4 con valore 4. \\
24 & Definisce la costante di compilazione M1\_\allowbreak{}DIRECTION\_\allowbreak{}CHANNEL con valore 10. \\
25 & Definisce una costante di PWM/frequenza/velocita' per M1\_\allowbreak{}SPEED\_\allowbreak{}CHANNEL. \\
26 & Definisce la costante di compilazione M2\_\allowbreak{}DIRECTION\_\allowbreak{}CHANNEL con valore 11. \\
27 & Definisce una costante di PWM/frequenza/velocita' per M2\_\allowbreak{}SPEED\_\allowbreak{}CHANNEL. \\
28 & Definisce la costante di compilazione M3\_\allowbreak{}DIRECTION\_\allowbreak{}CHANNEL con valore 4. \\
29 & Definisce una costante di PWM/frequenza/velocita' per M3\_\allowbreak{}SPEED\_\allowbreak{}CHANNEL. \\
30 & Definisce la costante di compilazione M4\_\allowbreak{}DIRECTION\_\allowbreak{}CHANNEL con valore 5. \\
31 & Definisce una costante di PWM/frequenza/velocita' per M4\_\allowbreak{}SPEED\_\allowbreak{}CHANNEL. \\
34 & Definisce la costante di compilazione FORWARD con valore 1. \\
35 & Definisce la costante di compilazione BACKWARD con valore 2. \\
38 & Definisce una costante di PWM/frequenza/velocita' per MAX\_\allowbreak{}SPEED. \\
41 & Definisce la costante di compilazione S1 con valore 0. \\
42 & Definisce la costante di compilazione S2 con valore 1. \\
43 & Definisce la costante di compilazione S3 con valore 3. \\
44 & Definisce la costante di compilazione S4 con valore 6. \\
45 & Definisce la costante di compilazione S5 con valore 9. \\
46 & Definisce la costante di compilazione S6 con valore 12. \\
47 & Definisce la costante di compilazione S7 con valore 14. \\
48 & Definisce la costante di compilazione S8 con valore 15. \\
50 & Dichiara la classe MotorController e separa l'interfaccia pubblica dai dettagli privati. \\
51 & Cambia la visibilita' dei membri della classe per il compilatore. \\
53 & Dichiara l'API MotorController senza implementarla: la definizione e' in un file .cpp o piu' avanti. \\
56 & Dichiara l'API setAddress senza implementarla: la definizione e' in un file .cpp o piu' avanti. \\
59 & Dichiara l'API begin senza implementarla: la definizione e' in un file .cpp o piu' avanti. \\
62 & Dichiara l'API DCrun senza implementarla: la definizione e' in un file .cpp o piu' avanti. \\
65 & Dichiara l'API DCbrake senza implementarla: la definizione e' in un file .cpp o piu' avanti. \\
66 & Dichiara l'API DCbrakeAll senza implementarla: la definizione e' in un file .cpp o piu' avanti. \\
67 & Dichiara l'API bipolarStepperStop senza implementarla: la definizione e' in un file .cpp o piu' avanti. \\
70 & Dichiara l'API servoTurn senza implementarla: la definizione e' in un file .cpp o piu' avanti. \\
73 & Dichiara l'API servoPairTurn senza implementarla: la definizione e' in un file .cpp o piu' avanti. \\
75 & Cambia la visibilita' dei membri della classe per il compilatore. \\
76 & Dichiara o inizializza una variabile: il tipo determina intervallo, durata e semantica del dato. \\
77 & Dichiara o inizializza una variabile: il tipo determina intervallo, durata e semantica del dato. \\
80 & Dichiara l'API setDigitalChannel senza implementarla: la definizione e' in un file .cpp o piu' avanti. \\
81 & Dichiara l'API setPwmChannel senza implementarla: la definizione e' in un file .cpp o piu' avanti. \\
84 & Chiude la include guard o il blocco di precompilazione aperto in precedenza. \\
\bottomrule
\end{longtable}


\clearpage
\subsection{Traduzione motori e servo in canali PCA}
\textbf{File:} \codepath{tank/lib/motorController/motorController.cpp}\\
\textbf{Collegamenti audit:} \codepath{A-02, A-11}

\subsubsection{Panoramica}
L'implementazione seleziona i due canali corretti per ogni motore, limita il PWM e scrive prima velocita' poi direzione. Per i servo converte gradi in tick PWM a 50 Hz. Tutte le operazioni finiscono nel \codepath{PWMController}, quindi la loro affidabilita' dipende dal bus I2C.

\subsubsection{Blocchi logici}
\begin{longtable}{p{0.14\textwidth}p{0.24\textwidth}p{0.50\textwidth}}
\toprule
\textbf{Righe} & \textbf{Blocco} & \textbf{Cosa studiare} \\
\midrule
\endhead
3--22 & Limiti servo e costruzione & Definisce tick servo 110--500, conserva frequenza e permette di cambiare indirizzo prima del begin. \\
24--54 & DCrun & Seleziona la coppia di canali con switch, limita velocita' e invia due scritture PCA. \\
56--90 & Stop motori & Spegne PWM e direzione per un motore o per tutti; il nome brake non sostituisce una verifica del comportamento H-bridge. \\
92--108 & Servo & Clampa angoli e calcola l'angolo specchiato della seconda uscita. \\
110--129 & Encoding PCA & Usa 4096 per i flag digitali e tre rami per off, full-on o PWM normale. \\
\bottomrule
\end{longtable}

\subsubsection{Funzioni e failure mode}
\codepath{DCrun} rifiuta un numero motore non riconosciuto. \codepath{DCbrakeAll} e' usato al boot. \codepath{servoPairTurn} impone un accoppiamento geometrico ma non conosce la meccanica reale. Non vengono restituiti errori I2C: perciò lo strato chiamante non sa se l'uscita ha davvero cambiato stato.

\subsubsection{Listing completo numerato}
\lstinputlisting[language=C++]{../tank/lib/motorController/motorController.cpp}

\subsubsection{Lettura riga per riga}
\begin{longtable}{p{0.10\textwidth}p{0.78\textwidth}}
\toprule
\textbf{Riga} & \textbf{Lettura riga per riga} \\
\midrule
\endhead
1 & Include il modulo locale motorController.h e rende visibili la sua interfaccia e i suoi tipi. \\
3 & Apre un namespace anonimo: costanti e helper restano privati a questo file di implementazione. \\
6 & Dichiara o inizializza una variabile: il tipo determina intervallo, durata e semantica del dato. \\
7 & Dichiara o inizializza una variabile: il tipo determina intervallo, durata e semantica del dato. \\
10 & Dichiara o inizializza una variabile: il tipo determina intervallo, durata e semantica del dato. \\
11 & Riga sintatticamente significativa del blocco corrente; il suo ruolo e' dettagliato dal contesto e dal listing adiacente. \\
13 & Definisce la funzione MotorController::setAddress; i parametri costituiscono il suo input e il corpo seguente ne realizza il contratto. \\
15 & Aggiorna o inizializza uno stato; seguire la variabile permette di vedere l'effetto sul ciclo successivo. \\
18 & Definisce la funzione MotorController::begin; i parametri costituiscono il suo input e il corpo seguente ne realizza il contratto. \\
20 & Esegue l'istruzione o la chiamata indicata nel contesto della funzione corrente. \\
21 & Esegue l'istruzione o la chiamata indicata nel contesto della funzione corrente. \\
24 & Definisce la funzione MotorController::DCrun; i parametri costituiscono il suo input e il corpo seguente ne realizza il contratto. \\
25 & Dichiara o inizializza una variabile: il tipo determina intervallo, durata e semantica del dato. \\
26 & Dichiara o inizializza una variabile: il tipo determina intervallo, durata e semantica del dato. \\
30 & Apre una selezione multi-ramo basata sul valore dell'espressione. \\
31 & Seleziona questo ramo dello switch quando il valore confrontato coincide con il caso indicato. \\
32 & Aggiorna o inizializza uno stato; seguire la variabile permette di vedere l'effetto sul ciclo successivo. \\
33 & Aggiorna o inizializza uno stato; seguire la variabile permette di vedere l'effetto sul ciclo successivo. \\
34 & Esce dal ramo corrente dello switch evitando il passaggio al caso successivo. \\
35 & Seleziona questo ramo dello switch quando il valore confrontato coincide con il caso indicato. \\
36 & Aggiorna o inizializza uno stato; seguire la variabile permette di vedere l'effetto sul ciclo successivo. \\
37 & Aggiorna o inizializza uno stato; seguire la variabile permette di vedere l'effetto sul ciclo successivo. \\
38 & Esce dal ramo corrente dello switch evitando il passaggio al caso successivo. \\
39 & Seleziona questo ramo dello switch quando il valore confrontato coincide con il caso indicato. \\
40 & Aggiorna o inizializza uno stato; seguire la variabile permette di vedere l'effetto sul ciclo successivo. \\
41 & Aggiorna o inizializza uno stato; seguire la variabile permette di vedere l'effetto sul ciclo successivo. \\
42 & Esce dal ramo corrente dello switch evitando il passaggio al caso successivo. \\
43 & Seleziona questo ramo dello switch quando il valore confrontato coincide con il caso indicato. \\
44 & Aggiorna o inizializza uno stato; seguire la variabile permette di vedere l'effetto sul ciclo successivo. \\
45 & Aggiorna o inizializza uno stato; seguire la variabile permette di vedere l'effetto sul ciclo successivo. \\
46 & Esce dal ramo corrente dello switch evitando il passaggio al caso successivo. \\
47 & Definisce il ramo di fallback dello switch per input non riconosciuti. \\
48 & Termina la funzione e restituisce al chiamante il valore o il controllo indicato. \\
52 & Compone un comando verso un canale PCA9685 o verso il relativo helper basso livello. \\
53 & Compone un comando verso un canale PCA9685 o verso il relativo helper basso livello. \\
56 & Definisce la funzione MotorController::DCbrake; i parametri costituiscono il suo input e il corpo seguente ne realizza il contratto. \\
58 & Apre una selezione multi-ramo basata sul valore dell'espressione. \\
59 & Seleziona questo ramo dello switch quando il valore confrontato coincide con il caso indicato. \\
60 & Compone un comando verso un canale PCA9685 o verso il relativo helper basso livello. \\
61 & Compone un comando verso un canale PCA9685 o verso il relativo helper basso livello. \\
62 & Esce dal ramo corrente dello switch evitando il passaggio al caso successivo. \\
63 & Seleziona questo ramo dello switch quando il valore confrontato coincide con il caso indicato. \\
64 & Compone un comando verso un canale PCA9685 o verso il relativo helper basso livello. \\
65 & Compone un comando verso un canale PCA9685 o verso il relativo helper basso livello. \\
66 & Esce dal ramo corrente dello switch evitando il passaggio al caso successivo. \\
67 & Seleziona questo ramo dello switch quando il valore confrontato coincide con il caso indicato. \\
68 & Compone un comando verso un canale PCA9685 o verso il relativo helper basso livello. \\
69 & Compone un comando verso un canale PCA9685 o verso il relativo helper basso livello. \\
70 & Esce dal ramo corrente dello switch evitando il passaggio al caso successivo. \\
71 & Seleziona questo ramo dello switch quando il valore confrontato coincide con il caso indicato. \\
72 & Compone un comando verso un canale PCA9685 o verso il relativo helper basso livello. \\
73 & Compone un comando verso un canale PCA9685 o verso il relativo helper basso livello. \\
74 & Esce dal ramo corrente dello switch evitando il passaggio al caso successivo. \\
78 & Definisce la funzione MotorController::DCbrakeAll; i parametri costituiscono il suo input e il corpo seguente ne realizza il contratto. \\
80 & Esegue l'istruzione o la chiamata indicata nel contesto della funzione corrente. \\
81 & Esegue l'istruzione o la chiamata indicata nel contesto della funzione corrente. \\
82 & Esegue l'istruzione o la chiamata indicata nel contesto della funzione corrente. \\
83 & Esegue l'istruzione o la chiamata indicata nel contesto della funzione corrente. \\
86 & Definisce la funzione MotorController::bipolarStepperStop; i parametri costituiscono il suo input e il corpo seguente ne realizza il contratto. \\
88 & Esegue l'istruzione o la chiamata indicata nel contesto della funzione corrente. \\
89 & Esegue l'istruzione o la chiamata indicata nel contesto della funzione corrente. \\
92 & Definisce la funzione MotorController::servoTurn; i parametri costituiscono il suo input e il corpo seguente ne realizza il contratto. \\
94 & Limita un valore all'intervallo ammesso prima dell'uso successivo. \\
95 & Rimappa linearmente un intervallo intero in un altro; il risultato usa aritmetica intera e puo' troncare. \\
96 & Compone un comando verso un canale PCA9685 o verso il relativo helper basso livello. \\
99 & Definisce la funzione MotorController::servoPairTurn; i parametri costituiscono il suo input e il corpo seguente ne realizza il contratto. \\
101 & Limita un valore all'intervallo ammesso prima dell'uso successivo. \\
104 & Limita un valore all'intervallo ammesso prima dell'uso successivo. \\
106 & Esegue l'istruzione o la chiamata indicata nel contesto della funzione corrente. \\
107 & Esegue l'istruzione o la chiamata indicata nel contesto della funzione corrente. \\
110 & Definisce la funzione MotorController::setDigitalChannel; i parametri costituiscono il suo input e il corpo seguente ne realizza il contratto. \\
113 & Esegue l'istruzione o la chiamata indicata nel contesto della funzione corrente. \\
116 & Definisce la funzione MotorController::setPwmChannel; i parametri costituiscono il suo input e il corpo seguente ne realizza il contratto. \\
118 & Valuta una guardia condizionale prima di eseguire il blocco successivo. \\
119 & Esegue l'istruzione o la chiamata indicata nel contesto della funzione corrente. \\
122 & Prosegue la catena condizionale: questo ramo e' valutato solo se quelli precedenti non sono stati scelti. \\
123 & Esegue l'istruzione o la chiamata indicata nel contesto della funzione corrente. \\
126 & Apre l'alternativa della condizione precedente. \\
127 & Esegue l'istruzione o la chiamata indicata nel contesto della funzione corrente. \\
\bottomrule
\end{longtable}


\clearpage
\subsection{Interfaccia minima PCA9685}
\textbf{File:} \codepath{tank/lib/PWMController/PWMController.h}\\
\textbf{Collegamenti audit:} \codepath{A-02, A-10}

\subsubsection{Panoramica}
Il driver basso livello conosce indirizzo I2C e registri PCA9685, ma non conosce motori o servo. La separazione consente di riutilizzare l'accesso a registri, ma impone agli strati superiori di controllare errori e validita' dei canali.

\subsubsection{Blocchi logici}
\begin{longtable}{p{0.14\textwidth}p{0.24\textwidth}p{0.50\textwidth}}
\toprule
\textbf{Righe} & \textbf{Blocco} & \textbf{Cosa studiare} \\
\midrule
\endhead
1--14 & Guardia e responsabilita' & Definisce lo scopo limitato: inizializzazione I2C, frequenza e registri PWM. \\
15--41 & Classe & Espone begin, frequenza e setPWM; conserva registri e indirizzo come dettagli privati. \\
\bottomrule
\end{longtable}

\subsubsection{Funzioni e failure mode}
\codepath{begin} inizializza il bus e i mode register. \codepath{setPWMFreq} programma il prescaler comune. \codepath{setPWM} scrive quattro byte di un canale. \codepath{readRegister} e \codepath{writeRegister} sono primitive I2C private.

\subsubsection{Listing completo numerato}
\lstinputlisting[language=C++]{../tank/lib/PWMController/PWMController.h}

\subsubsection{Lettura riga per riga}
\begin{longtable}{p{0.10\textwidth}p{0.78\textwidth}}
\toprule
\textbf{Riga} & \textbf{Lettura riga per riga} \\
\midrule
\endhead
1 & Inizia una include guard: il corpo dell'header verra' elaborato solo se la macro non e' gia' definita. \\
2 & Definisce la macro della include guard: impedisce inclusioni multiple dell'header. \\
4 & Importa il core Arduino: tipi, GPIO, tempo, map e constrain. \\
15 & Dichiara la classe PWMController e separa l'interfaccia pubblica dai dettagli privati. \\
16 & Cambia la visibilita' dei membri della classe per il compilatore. \\
17 & Dichiara l'API PWMController senza implementarla: la definizione e' in un file .cpp o piu' avanti. \\
20 & Dichiara l'API begin senza implementarla: la definizione e' in un file .cpp o piu' avanti. \\
23 & Dichiara l'API setPWMFreq senza implementarla: la definizione e' in un file .cpp o piu' avanti. \\
26 & Dichiara l'API setPWM senza implementarla: la definizione e' in un file .cpp o piu' avanti. \\
28 & Cambia la visibilita' dei membri della classe per il compilatore. \\
30 & Dichiara stato statico: dura per l'intera esecuzione e, se locale a funzione, conserva valore fra chiamate. \\
31 & Dichiara stato statico: dura per l'intera esecuzione e, se locale a funzione, conserva valore fra chiamate. \\
32 & Dichiara stato statico: dura per l'intera esecuzione e, se locale a funzione, conserva valore fra chiamate. \\
33 & Dichiara stato statico: dura per l'intera esecuzione e, se locale a funzione, conserva valore fra chiamate. \\
36 & Dichiara o inizializza una variabile: il tipo determina intervallo, durata e semantica del dato. \\
39 & Dichiara l'API readRegister senza implementarla: la definizione e' in un file .cpp o piu' avanti. \\
40 & Dichiara l'API writeRegister senza implementarla: la definizione e' in un file .cpp o piu' avanti. \\
43 & Chiude la include guard o il blocco di precompilazione aperto in precedenza. \\
\bottomrule
\end{longtable}


\clearpage
\subsection{Scritture I2C verso PCA9685}
\textbf{File:} \codepath{tank/lib/PWMController/PWMController.cpp}\\
\textbf{Collegamenti audit:} \codepath{A-02}

\subsubsection{Panoramica}
Questo e' il punto piu' vicino alla periferica PWM. A 50 Hz il prescaler calcolato e' circa 124. Ogni canale richiede registro base piu' quattro byte ON/OFF. Il codice non controlla il risultato delle transazioni, quindi un errore I2C non raggiunge il controller dei cingoli.

\subsubsection{Blocchi logici}
\begin{longtable}{p{0.14\textwidth}p{0.24\textwidth}p{0.50\textwidth}}
\toprule
\textbf{Righe} & \textbf{Blocco} & \textbf{Cosa studiare} \\
\midrule
\endhead
1--17 & Costruzione e begin & Memorizza indirizzo, avvia Wire e programma MODE1/MODE2. \\
19--36 & Frequenza & Corregge il clock, calcola/arrotonda prescaler, entra in sleep, ripristina mode e aspetta 5 ms. \\
38--48 & Scrittura PWM & Indirizza il primo registro del canale e invia i quattro byte in little-endian logico. \\
50--67 & Lettura e scrittura registro & Esegue transazioni minimali; in caso di assenza dati la lettura restituisce zero. \\
\bottomrule
\end{longtable}

\subsubsection{Funzioni e failure mode}
\codepath{setPWMFreq} usa la formula del datasheet. \codepath{setPWM} assume che il canale sia 0--15 e che la transazione vada a buon fine. \codepath{readRegister} interpreta nessun byte come zero, valore che puo' confondersi con un registro realmente zero. \codepath{writeRegister} non propaga l'esito di \codepath{Wire.endTransmission}.

\subsubsection{Listing completo numerato}
\lstinputlisting[language=C++]{../tank/lib/PWMController/PWMController.cpp}

\subsubsection{Lettura riga per riga}
\begin{longtable}{p{0.10\textwidth}p{0.78\textwidth}}
\toprule
\textbf{Riga} & \textbf{Lettura riga per riga} \\
\midrule
\endhead
1 & Include il modulo locale PWMController.h e rende visibili la sua interfaccia e i suoi tipi. \\
3 & Importa l'API I2C Wire usata per parlare con la PCA9685. \\
5 & Dichiara o inizializza una variabile: il tipo determina intervallo, durata e semantica del dato. \\
7 & Definisce la funzione PWMController::begin; i parametri costituiscono il suo input e il corpo seguente ne realizza il contratto. \\
9 & Esegue un'operazione I2C sul bus condiviso; il suo esito e la sua durata influenzano il loop cooperativo. \\
12 & Esegue l'istruzione o la chiamata indicata nel contesto della funzione corrente. \\
16 & Esegue l'istruzione o la chiamata indicata nel contesto della funzione corrente. \\
19 & Definisce la funzione PWMController::setPWMFreq; i parametri costituiscono il suo input e il corpo seguente ne realizza il contratto. \\
21 & Dichiara o inizializza una variabile: il tipo determina intervallo, durata e semantica del dato. \\
24 & Dichiara o inizializza una variabile: il tipo determina intervallo, durata e semantica del dato. \\
25 & Dichiara o inizializza una variabile: il tipo determina intervallo, durata e semantica del dato. \\
27 & Dichiara o inizializza una variabile: il tipo determina intervallo, durata e semantica del dato. \\
28 & Dichiara o inizializza una variabile: il tipo determina intervallo, durata e semantica del dato. \\
31 & Esegue l'istruzione o la chiamata indicata nel contesto della funzione corrente. \\
32 & Esegue l'istruzione o la chiamata indicata nel contesto della funzione corrente. \\
33 & Esegue l'istruzione o la chiamata indicata nel contesto della funzione corrente. \\
34 & Blocca intenzionalmente il loop per il tempo indicato; non e' un timer concorrente. \\
35 & Esegue l'istruzione o la chiamata indicata nel contesto della funzione corrente. \\
38 & Definisce la funzione PWMController::setPWM; i parametri costituiscono il suo input e il corpo seguente ne realizza il contratto. \\
41 & Esegue un'operazione I2C sul bus condiviso; il suo esito e la sua durata influenzano il loop cooperativo. \\
42 & Esegue un'operazione I2C sul bus condiviso; il suo esito e la sua durata influenzano il loop cooperativo. \\
43 & Esegue un'operazione I2C sul bus condiviso; il suo esito e la sua durata influenzano il loop cooperativo. \\
44 & Esegue un'operazione I2C sul bus condiviso; il suo esito e la sua durata influenzano il loop cooperativo. \\
45 & Esegue un'operazione I2C sul bus condiviso; il suo esito e la sua durata influenzano il loop cooperativo. \\
46 & Esegue un'operazione I2C sul bus condiviso; il suo esito e la sua durata influenzano il loop cooperativo. \\
47 & Esegue un'operazione I2C sul bus condiviso; il suo esito e la sua durata influenzano il loop cooperativo. \\
50 & Definisce la funzione PWMController::readRegister; i parametri costituiscono il suo input e il corpo seguente ne realizza il contratto. \\
52 & Esegue un'operazione I2C sul bus condiviso; il suo esito e la sua durata influenzano il loop cooperativo. \\
53 & Esegue un'operazione I2C sul bus condiviso; il suo esito e la sua durata influenzano il loop cooperativo. \\
54 & Esegue un'operazione I2C sul bus condiviso; il suo esito e la sua durata influenzano il loop cooperativo. \\
57 & Esegue un'operazione I2C sul bus condiviso; il suo esito e la sua durata influenzano il loop cooperativo. \\
58 & Termina la funzione e restituisce al chiamante il valore o il controllo indicato. \\
61 & Definisce la funzione PWMController::writeRegister; i parametri costituiscono il suo input e il corpo seguente ne realizza il contratto. \\
63 & Esegue un'operazione I2C sul bus condiviso; il suo esito e la sua durata influenzano il loop cooperativo. \\
64 & Esegue un'operazione I2C sul bus condiviso; il suo esito e la sua durata influenzano il loop cooperativo. \\
65 & Esegue un'operazione I2C sul bus condiviso; il suo esito e la sua durata influenzano il loop cooperativo. \\
66 & Esegue un'operazione I2C sul bus condiviso; il suo esito e la sua durata influenzano il loop cooperativo. \\
\bottomrule
\end{longtable}


\clearpage
\subsection{Configurazione PlatformIO del controller}
\textbf{File:} \codepath{controller-esp32/platformio.ini}\\
\textbf{Collegamenti audit:} \codepath{A-10}

\subsubsection{Panoramica}
Seleziona un ESP32 Dev Module generico, framework Arduino, monitor a 115200 baud e velocita' di upload. \codepath{monitor_speed} coincide con \codepath{Serial.begin(115200)}. Come il tank, non fissa una versione della platform: il risultato futuro puo' dipendere dalla risoluzione delle dipendenze.

\subsubsection{Blocchi logici}
\begin{longtable}{p{0.14\textwidth}p{0.24\textwidth}p{0.50\textwidth}}
\toprule
\textbf{Righe} & \textbf{Blocco} & \textbf{Cosa studiare} \\
\midrule
\endhead
2--7 & Ambiente ESP32 & Indica board, platform, framework, seriale e upload. \\
\bottomrule
\end{longtable}

\subsubsection{Funzioni e failure mode}
Non contiene funzioni. E' letto da PlatformIO e determina toolchain e API disponibili al firmware.

\subsubsection{Listing completo numerato}
\lstinputlisting[language=]{../controller-esp32/platformio.ini}

\subsubsection{Lettura riga per riga}
\begin{longtable}{p{0.10\textwidth}p{0.78\textwidth}}
\toprule
\textbf{Riga} & \textbf{Lettura riga per riga} \\
\midrule
\endhead
2 & Riga sintatticamente significativa del blocco corrente; il suo ruolo e' dettagliato dal contesto e dal listing adiacente. \\
3 & Aggiorna o inizializza uno stato; seguire la variabile permette di vedere l'effetto sul ciclo successivo. \\
4 & Aggiorna o inizializza uno stato; seguire la variabile permette di vedere l'effetto sul ciclo successivo. \\
5 & Aggiorna o inizializza uno stato; seguire la variabile permette di vedere l'effetto sul ciclo successivo. \\
6 & Aggiorna o inizializza uno stato; seguire la variabile permette di vedere l'effetto sul ciclo successivo. \\
7 & Aggiorna o inizializza uno stato; seguire la variabile permette di vedere l'effetto sul ciclo successivo. \\
\bottomrule
\end{longtable}


\clearpage
\subsection{Scheduler del controller ESP32}
\textbf{File:} \codepath{controller-esp32/src/main.cpp}\\
\textbf{Collegamenti audit:} \codepath{A-07, A-09, A-11}

\subsubsection{Panoramica}
Il controller usa un loop cooperativo minimo: ogni almeno 10 ms acquisisce input e prova a inviare un datagramma. Non attende il completamento WiFi nel setup. Il periodo e' un limite minimo, non una garanzia esatta: ADC, rete e seriale possono aggiungere tempo.

\subsubsection{Blocchi logici}
\begin{longtable}{p{0.14\textwidth}p{0.24\textwidth}p{0.50\textwidth}}
\toprule
\textbf{Righe} & \textbf{Blocco} & \textbf{Cosa studiare} \\
\midrule
\endhead
1--20 & Include, intervallo e stato & Importa i moduli, definisce 10 ms e conserva il timestamp dell'ultimo invio. \\
22--30 & setup & Avvia seriale, calibra/configura ingressi e avvia la procedura WiFi. \\
32--42 & loop & Ritorna prima dello scadere del periodo, poi legge ControllerData e invia UDP. \\
\bottomrule
\end{longtable}

\subsubsection{Funzioni e failure mode}
\codepath{setup} crea le condizioni iniziali. \codepath{loop} e' uno scheduler a polling: non usa task, timer hardware o ISR del progetto. La sottrazione unsigned con \codepath{millis} resta corretta al rollover, ma non recupera gli intervalli persi.

\subsubsection{Listing completo numerato}
\lstinputlisting[language=C++]{../controller-esp32/src/main.cpp}

\subsubsection{Lettura riga per riga}
\begin{longtable}{p{0.10\textwidth}p{0.78\textwidth}}
\toprule
\textbf{Riga} & \textbf{Lettura riga per riga} \\
\midrule
\endhead
1 & Importa il core Arduino: tipi, GPIO, tempo, map e constrain. \\
3 & Include il modulo locale joystickReader.h e rende visibili la sua interfaccia e i suoi tipi. \\
4 & Include il modulo locale udpSender.h e rende visibili la sua interfaccia e i suoi tipi. \\
18 & Definisce una durata nominale in millisecondi per UDP\_\allowbreak{}SEND\_\allowbreak{}INTERVAL\_\allowbreak{}MS; viene rispettata solo quando il loop continua a girare. \\
20 & Dichiara stato statico: dura per l'intera esecuzione e, se locale a funzione, conserva valore fra chiamate. \\
22 & Definisce la funzione setup; i parametri costituiscono il suo input e il corpo seguente ne realizza il contratto. \\
23 & Emette diagnostica seriale; e' osservabilita' utile ma puo' aggiungere tempo al giro di loop. \\
26 & Esegue l'istruzione o la chiamata indicata nel contesto della funzione corrente. \\
27 & Esegue l'istruzione o la chiamata indicata nel contesto della funzione corrente. \\
29 & Emette diagnostica seriale; e' osservabilita' utile ma puo' aggiungere tempo al giro di loop. \\
32 & Definisce la funzione loop; i parametri costituiscono il suo input e il corpo seguente ne realizza il contratto. \\
33 & Legge o confronta il contatore temporale Arduino; la sottrazione unsigned e' rollover-safe ma richiede un loop che torni a eseguire. \\
34 & Valuta una guardia condizionale prima di eseguire il blocco successivo. \\
35 & Termina la funzione e restituisce al chiamante il valore o il controllo indicato. \\
37 & Aggiorna o inizializza uno stato; seguire la variabile permette di vedere l'effetto sul ciclo successivo. \\
40 & Dichiara o inizializza una variabile: il tipo determina intervallo, durata e semantica del dato. \\
41 & Esegue l'istruzione o la chiamata indicata nel contesto della funzione corrente. \\
\bottomrule
\end{longtable}


\clearpage
\subsection{Contratto dei dati del controller}
\textbf{File:} \codepath{controller-esp32/src/joystickReader.h}\\
\textbf{Collegamenti audit:} \codepath{A-05, A-07, A-12, A-14}

\subsubsection{Panoramica}
\codepath{ControllerData} e' il contratto fra lettura hardware e trasporto UDP. I quattro assi restano \codepath{int} in scala 0--1023; i pulsanti sono booleani. La struct evita che il sender debba conoscere i pin o la calibrazione.

\subsubsection{Blocchi logici}
\begin{longtable}{p{0.14\textwidth}p{0.24\textwidth}p{0.50\textwidth}}
\toprule
\textbf{Righe} & \textbf{Blocco} & \textbf{Cosa studiare} \\
\midrule
\endhead
1--4 & Guardia e Arduino & Protegge l'header e rende disponibili tipi/API Arduino. \\
14--27 & ControllerData & Raggruppa assi e pulsanti con nomi coerenti con il protocollo V1. \\
29--37 & API & Dichiara inizializzazione, riarmo neutro dopo reconnect e lettura completa di un campione. \\
\bottomrule
\end{longtable}

\subsubsection{Funzioni e failure mode}
\codepath{JoystickReader_begin} configura hardware e calibra. \codepath{JoystickReader_requireNeutralBeforeCommands} riporta il modulo nello stato non armato dopo boot/reconnect. \codepath{JoystickReader_read} restituisce sempre una nuova struct, pronta per la serializzazione oppure sicura quando calibrazione o riarmo non sono validi.

\subsubsection{Listing completo numerato}
\lstinputlisting[language=C++]{../controller-esp32/src/joystickReader.h}

\subsubsection{Lettura riga per riga}
\begin{longtable}{p{0.10\textwidth}p{0.78\textwidth}}
\toprule
\textbf{Riga} & \textbf{Lettura riga per riga} \\
\midrule
\endhead
1 & Inizia una include guard: il corpo dell'header verra' elaborato solo se la macro non e' gia' definita. \\
2 & Definisce la macro della include guard: impedisce inclusioni multiple dell'header. \\
4 & Importa il core Arduino: tipi, GPIO, tempo, map e constrain. \\
15 & Dichiara la struct ControllerData: un contenitore di campi che viaggiano insieme fra moduli. \\
17 & Dichiara o inizializza una variabile: il tipo determina intervallo, durata e semantica del dato. \\
18 & Dichiara o inizializza una variabile: il tipo determina intervallo, durata e semantica del dato. \\
21 & Dichiara o inizializza una variabile: il tipo determina intervallo, durata e semantica del dato. \\
22 & Dichiara o inizializza una variabile: il tipo determina intervallo, durata e semantica del dato. \\
25 & Dichiara o inizializza una variabile: il tipo determina intervallo, durata e semantica del dato. \\
26 & Dichiara o inizializza una variabile: il tipo determina intervallo, durata e semantica del dato. \\
30 & Dichiara \codepath{JoystickReader_begin()}, che configura pulsanti/ADC, calibra e richiede neutralita' iniziale. \\
34 & Dichiara \codepath{JoystickReader_requireNeutralBeforeCommands()}, invocata anche dal sender quando la connessione UDP viene ristabilita. \\
37 & Dichiara \codepath{JoystickReader_read()}, che restituisce dati normalizzati oppure il campione sicuro. \\
39 & Chiude la include guard o il blocco di precompilazione aperto in precedenza. \\
\bottomrule
\end{longtable}


\clearpage
\subsection{ADC, calibrazione e pulsanti ESP32}
\textbf{File:} \codepath{controller-esp32/src/joystickReader.cpp}\\
\textbf{Collegamenti audit:} \codepath{A-05, A-07, A-12, A-14}

\subsubsection{Panoramica}
Il file legge quattro ingressi ADC1, adatta i 12 bit ESP32 alla scala 0--1023 e corregge il centro dei joystick. GPIO32--35 appartengono ad ADC1 e sono quindi adatti a WiFi sul classico ESP32; GPIO25/26 usano \codepath{INPUT_PULLUP}, per cui LOW significa pulsante premuto. La guida Joy1 applica \codepath{DRIVE_INPUT_DEADZONE=20}; Joy2 mantiene \codepath{TURRET_INPUT_DEADZONE=80}. Sono filtri di rumore all'ingresso, non la protezione motore del tank. La calibrazione blocca intenzionalmente l'avvio per almeno circa 1,24 s: 600 ms iniziali piu' quattro serie da 80 campioni con delay da 2 ms.

\subsubsection{Blocchi logici}
\begin{longtable}{p{0.14\textwidth}p{0.24\textwidth}p{0.50\textwidth}}
\toprule
\textbf{Righe} & \textbf{Blocco} & \textbf{Cosa studiare} \\
\midrule
\endhead
15--76 & Pin, scala, filtri e stati & Definisce ingressi, centro, range, deadzone separate, swap/inversioni e stato di calibrazione/riarmo. \\
78--86 & Stato pulsanti & Dichiara il debounce con candidato, valore stabile e timestamp per Zero e Fire. \\
88--141 & ADC, calibrazione asse e ricentraggio & Legge ADC, rimappa, valida centro/spread e corregge ciascun semiasse con la deadzone ricevuta come argomento. \\
143--151 & Inversione asse & Limita e inverte un asse normalizzato, preservando esattamente il centro. \\
153--185 & Procedura calibrazione & Mostra messaggi, aspetta 600 ms, calibra i quattro assi e restituisce un solo esito booleano. \\
187--219 & Debounce e campione sicuro & Inizializza/legge pulsanti, costruisce dati sicuri e verifica i quattro assi neutrali. \\
221--238 & Begin e riarmo & Configura pull-up/ADC, calibra e invalida l'abilitazione dei comandi. \\
240--297 & Lettura completa & Acquisisce assi con deadzone corretta, gestisce riarmo, swap/inversioni e restituisce il frame locale. \\
\bottomrule
\end{longtable}

\subsubsection{Funzioni e failure mode}
\codepath{readRawJoystick} converte 0--4095 in 0--1023. \codepath{calibrateAxisCenter} calcola una media sicura: $80 \times 1023 = 81840$, entro \codepath{long}, e rifiuta assi troppo lontani dal centro o instabili. \codepath{readCalibratedJoystick} ricostruisce i due semiassi rispetto al centro misurato; riceve la deadzone come parametro, cosi Joy1 usa 20 e Joy2 80 senza duplicare la funzione. \codepath{invertAxisIfNeeded} usa 1023-valore, preservando esplicitamente 512. \codepath{calibrateJoysticks} blocca intenzionalmente il boot. \codepath{readDebouncedButton}, \codepath{safeControllerData} e \codepath{axesAreNeutral} sostengono il riarmo. \codepath{JoystickReader_read} realizza il mapping finale: driveX=GPIO32, driveY=GPIO34, turretX=GPIO33, elevationY=GPIO35; le inversioni di Joy1 vengono applicate qui e non sul tank.

\subsubsection{Listing completo numerato}
\lstinputlisting[language=C++]{../controller-esp32/src/joystickReader.cpp}

\subsubsection{Lettura riga per riga}
\begin{longtable}{p{0.10\textwidth}p{0.78\textwidth}}
\toprule
\textbf{Riga} & \textbf{Lettura riga per riga} \\
\midrule
\endhead
1 & Include il modulo locale joystickReader.h e rende visibili la sua interfaccia e i suoi tipi. \\
16 & Configura il pin hardware tramite la macro DRIVE\_\allowbreak{}JOYSTICK\_\allowbreak{}X\_\allowbreak{}PIN; il valore associato deve coincidere con il cablaggio. \\
17 & Configura il pin hardware tramite la macro DRIVE\_\allowbreak{}JOYSTICK\_\allowbreak{}Y\_\allowbreak{}PIN; il valore associato deve coincidere con il cablaggio. \\
20 & Configura il pin hardware tramite la macro TURRET\_\allowbreak{}JOYSTICK\_\allowbreak{}X\_\allowbreak{}PIN; il valore associato deve coincidere con il cablaggio. \\
21 & Configura il pin hardware tramite la macro TURRET\_\allowbreak{}JOYSTICK\_\allowbreak{}Y\_\allowbreak{}PIN; il valore associato deve coincidere con il cablaggio. \\
24 & Configura il pin hardware tramite la macro ZERO\_\allowbreak{}BUTTON\_\allowbreak{}PIN; il valore associato deve coincidere con il cablaggio. \\
25 & Configura il pin hardware tramite la macro FIRE\_\allowbreak{}BUTTON\_\allowbreak{}PIN; il valore associato deve coincidere con il cablaggio. \\
28 & Definisce una soglia di neutralita' del joystick per JOYSTICK\_\allowbreak{}CENTER. \\
29 & Definisce la costante di compilazione AXIS\_\allowbreak{}MIN con valore 0. \\
30 & Definisce la costante di compilazione AXIS\_\allowbreak{}MAX con valore 1023. \\
33 & Definisce una soglia di neutralita' del joystick per CONTROLLER\_\allowbreak{}DEADZONE. \\
36 & Definisce la costante di compilazione CALIBRATION\_\allowbreak{}SAMPLES con valore 80. \\
39 & Configura il verso o lo scambio logico dell'asse DRIVE\_\allowbreak{}SWAP\_\allowbreak{}X\_\allowbreak{}Y. \\
40 & Configura il verso o lo scambio logico dell'asse TURRET\_\allowbreak{}SWAP\_\allowbreak{}X\_\allowbreak{}Y. \\
44 & Configura il verso o lo scambio logico dell'asse DRIVE\_\allowbreak{}X\_\allowbreak{}INVERTED. \\
45 & Configura il verso o lo scambio logico dell'asse DRIVE\_\allowbreak{}Y\_\allowbreak{}INVERTED. \\
46 & Configura il verso o lo scambio logico dell'asse TURRET\_\allowbreak{}X\_\allowbreak{}INVERTED. \\
47 & Configura il verso o lo scambio logico dell'asse ELEVATION\_\allowbreak{}Y\_\allowbreak{}INVERTED. \\
49 & Dichiara stato statico: dura per l'intera esecuzione e, se locale a funzione, conserva valore fra chiamate. \\
50 & Dichiara stato statico: dura per l'intera esecuzione e, se locale a funzione, conserva valore fra chiamate. \\
51 & Dichiara stato statico: dura per l'intera esecuzione e, se locale a funzione, conserva valore fra chiamate. \\
52 & Dichiara stato statico: dura per l'intera esecuzione e, se locale a funzione, conserva valore fra chiamate. \\
54 & Definisce la funzione readRawJoystick; i parametri costituiscono il suo input e il corpo seguente ne realizza il contratto. \\
55 & Acquisisce un campione ADC; il risultato dipende da tensione, attenuazione, rumore e cablaggio. \\
58 & Termina la funzione e restituisce al chiamante il valore o il controllo indicato. \\
61 & Definisce la funzione calibrateAxisCenter; i parametri costituiscono il suo input e il corpo seguente ne realizza il contratto. \\
62 & Dichiara o inizializza una variabile: il tipo determina intervallo, durata e semantica del dato. \\
64 & Avvia un ciclo con inizializzazione, condizione di prosecuzione e avanzamento espliciti. \\
65 & Aggiorna o inizializza uno stato; seguire la variabile permette di vedere l'effetto sul ciclo successivo. \\
66 & Blocca intenzionalmente il loop per il tempo indicato; non e' un timer concorrente. \\
69 & Termina la funzione e restituisce al chiamante il valore o il controllo indicato. \\
72 & Definisce la funzione readCalibratedJoystick; i parametri costituiscono il suo input e il corpo seguente ne realizza il contratto. \\
73 & Dichiara o inizializza una variabile: il tipo determina intervallo, durata e semantica del dato. \\
76 & Valuta una guardia condizionale prima di eseguire il blocco successivo. \\
77 & Termina la funzione e restituisce al chiamante il valore o il controllo indicato. \\
83 & Valuta una guardia condizionale prima di eseguire il blocco successivo. \\
84 & Dichiara o inizializza una variabile: il tipo determina intervallo, durata e semantica del dato. \\
85 & Termina la funzione e restituisce al chiamante il valore o il controllo indicato. \\
86 & Esegue l'istruzione o la chiamata indicata nel contesto della funzione corrente. \\
89 & Dichiara o inizializza una variabile: il tipo determina intervallo, durata e semantica del dato. \\
90 & Termina la funzione e restituisce al chiamante il valore o il controllo indicato. \\
91 & Esegue l'istruzione o la chiamata indicata nel contesto della funzione corrente. \\
94 & Definisce la funzione invertAxisIfNeeded; i parametri costituiscono il suo input e il corpo seguente ne realizza il contratto. \\
95 & Limita un valore all'intervallo ammesso prima dell'uso successivo. \\
97 & Valuta una guardia condizionale prima di eseguire il blocco successivo. \\
98 & Termina la funzione e restituisce al chiamante il valore o il controllo indicato. \\
101 & Termina la funzione e restituisce al chiamante il valore o il controllo indicato. \\
104 & Definisce la funzione calibrateJoysticks; i parametri costituiscono il suo input e il corpo seguente ne realizza il contratto. \\
105 & Emette diagnostica seriale; e' osservabilita' utile ma puo' aggiungere tempo al giro di loop. \\
106 & Emette diagnostica seriale; e' osservabilita' utile ma puo' aggiungere tempo al giro di loop. \\
107 & Emette diagnostica seriale; e' osservabilita' utile ma puo' aggiungere tempo al giro di loop. \\
108 & Blocca intenzionalmente il loop per il tempo indicato; non e' un timer concorrente. \\
110 & Aggiorna o inizializza uno stato; seguire la variabile permette di vedere l'effetto sul ciclo successivo. \\
111 & Aggiorna o inizializza uno stato; seguire la variabile permette di vedere l'effetto sul ciclo successivo. \\
112 & Aggiorna o inizializza uno stato; seguire la variabile permette di vedere l'effetto sul ciclo successivo. \\
113 & Aggiorna o inizializza uno stato; seguire la variabile permette di vedere l'effetto sul ciclo successivo. \\
115 & Emette diagnostica seriale; e' osservabilita' utile ma puo' aggiungere tempo al giro di loop. \\
116 & Emette diagnostica seriale; e' osservabilita' utile ma puo' aggiungere tempo al giro di loop. \\
117 & Emette diagnostica seriale; e' osservabilita' utile ma puo' aggiungere tempo al giro di loop. \\
118 & Emette diagnostica seriale; e' osservabilita' utile ma puo' aggiungere tempo al giro di loop. \\
119 & Emette diagnostica seriale; e' osservabilita' utile ma puo' aggiungere tempo al giro di loop. \\
121 & Emette diagnostica seriale; e' osservabilita' utile ma puo' aggiungere tempo al giro di loop. \\
122 & Emette diagnostica seriale; e' osservabilita' utile ma puo' aggiungere tempo al giro di loop. \\
123 & Emette diagnostica seriale; e' osservabilita' utile ma puo' aggiungere tempo al giro di loop. \\
124 & Emette diagnostica seriale; e' osservabilita' utile ma puo' aggiungere tempo al giro di loop. \\
125 & Emette diagnostica seriale; e' osservabilita' utile ma puo' aggiungere tempo al giro di loop. \\
126 & Emette diagnostica seriale; e' osservabilita' utile ma puo' aggiungere tempo al giro di loop. \\
127 & Emette diagnostica seriale; e' osservabilita' utile ma puo' aggiungere tempo al giro di loop. \\
130 & Definisce la funzione JoystickReader\_\allowbreak{}begin; i parametri costituiscono il suo input e il corpo seguente ne realizza il contratto. \\
132 & Configura elettricamente la direzione del GPIO prima di leggerlo o pilotarlo. \\
133 & Configura elettricamente la direzione del GPIO prima di leggerlo o pilotarlo. \\
135 & Configura il convertitore ADC ESP32 prima della lettura degli assi. \\
136 & Configura il convertitore ADC ESP32 prima della lettura degli assi. \\
138 & Esegue l'istruzione o la chiamata indicata nel contesto della funzione corrente. \\
141 & Definisce la funzione JoystickReader\_\allowbreak{}read; i parametri costituiscono il suo input e il corpo seguente ne realizza il contratto. \\
142 & Dichiara o inizializza una variabile: il tipo determina intervallo, durata e semantica del dato. \\
144 & Dichiara o inizializza una variabile: il tipo determina intervallo, durata e semantica del dato. \\
145 & Dichiara o inizializza una variabile: il tipo determina intervallo, durata e semantica del dato. \\
146 & Dichiara o inizializza una variabile: il tipo determina intervallo, durata e semantica del dato. \\
147 & Dichiara o inizializza una variabile: il tipo determina intervallo, durata e semantica del dato. \\
150 & Dichiara o inizializza una variabile: il tipo determina intervallo, durata e semantica del dato. \\
151 & Dichiara o inizializza una variabile: il tipo determina intervallo, durata e semantica del dato. \\
152 & Aggiorna o inizializza uno stato; seguire la variabile permette di vedere l'effetto sul ciclo successivo. \\
153 & Aggiorna o inizializza uno stato; seguire la variabile permette di vedere l'effetto sul ciclo successivo. \\
156 & Dichiara o inizializza una variabile: il tipo determina intervallo, durata e semantica del dato. \\
157 & Dichiara o inizializza una variabile: il tipo determina intervallo, durata e semantica del dato. \\
158 & Aggiorna o inizializza uno stato; seguire la variabile permette di vedere l'effetto sul ciclo successivo. \\
159 & Aggiorna o inizializza uno stato; seguire la variabile permette di vedere l'effetto sul ciclo successivo. \\
162 & Aggiorna o inizializza uno stato; seguire la variabile permette di vedere l'effetto sul ciclo successivo. \\
163 & Aggiorna o inizializza uno stato; seguire la variabile permette di vedere l'effetto sul ciclo successivo. \\
165 & Termina la funzione e restituisce al chiamante il valore o il controllo indicato. \\
\bottomrule
\end{longtable}


\clearpage
\subsection{Interfaccia del sender UDP}
\textbf{File:} \codepath{controller-esp32/src/udpSender.h}\\
\textbf{Collegamenti audit:} \codepath{A-04, A-09, A-10}

\subsubsection{Panoramica}
L'header dichiara il modulo che collega \codepath{ControllerData} a WiFi/UDP. Il tipo di ritorno \codepath{void} di \codepath{UdpSender_send} indica che il chiamante non riceve una conferma di consegna dal tank.

\subsubsection{Blocchi logici}
\begin{longtable}{p{0.14\textwidth}p{0.24\textwidth}p{0.50\textwidth}}
\toprule
\textbf{Righe} & \textbf{Blocco} & \textbf{Cosa studiare} \\
\midrule
\endhead
1--6 & Guardia e dipendenze & Include Arduino e il contratto dati del joystick. \\
16--22 & API & Dichiara avvio rete e invio del frame V1. \\
\bottomrule
\end{longtable}

\subsubsection{Funzioni e failure mode}
\codepath{UdpSender_begin} avvia la connessione station. \codepath{UdpSender_send} tenta l'invio solo quando WiFi risulta associato, ma l'associazione non e' conferma di ricezione UDP dal tank.

\subsubsection{Listing completo numerato}
\lstinputlisting[language=C++]{../controller-esp32/src/udpSender.h}

\subsubsection{Lettura riga per riga}
\begin{longtable}{p{0.10\textwidth}p{0.78\textwidth}}
\toprule
\textbf{Riga} & \textbf{Lettura riga per riga} \\
\midrule
\endhead
1 & Inizia una include guard: il corpo dell'header verra' elaborato solo se la macro non e' gia' definita. \\
2 & Definisce la macro della include guard: impedisce inclusioni multiple dell'header. \\
4 & Importa il core Arduino: tipi, GPIO, tempo, map e constrain. \\
6 & Include il modulo locale joystickReader.h e rende visibili la sua interfaccia e i suoi tipi. \\
17 & Dichiara l'API UdpSender\_\allowbreak{}begin senza implementarla: la definizione e' in un file .cpp o piu' avanti. \\
20 & Dichiara l'API UdpSender\_\allowbreak{}send senza implementarla: la definizione e' in un file .cpp o piu' avanti. \\
22 & Chiude la include guard o il blocco di precompilazione aperto in precedenza. \\
\bottomrule
\end{longtable}


\clearpage
\subsection{Connessione station e serializzazione UDP}
\textbf{File:} \codepath{controller-esp32/src/udpSender.cpp}\\
\textbf{Collegamenti audit:} \codepath{A-04, A-05, A-09, A-10}

\subsubsection{Panoramica}
Il controller opera come station della rete creata dal tank. Conserva socket, IP di fallback, timer e flag di transizione. Quando WiFi torna connesso usa il gateway come indirizzo del tank e riapre il socket. Il messaggio V1 e' costruito con \codepath{snprintf}; il frame normale massimo \codepath{V1;1023;1023;1023;1023;1;1} e' lungo 26 byte, quindi entra in \codepath{message[64]}.

\subsubsection{Blocchi logici}
\begin{longtable}{p{0.14\textwidth}p{0.24\textwidth}p{0.50\textwidth}}
\toprule
\textbf{Righe} & \textbf{Blocco} & \textbf{Cosa studiare} \\
\midrule
\endhead
1--39 & Librerie, parametri e stato & Definisce credenziali, porta, periodi, socket, IP e timer persistenti. \\
41--58 & Configurazione socket & Legge gateway DHCP, applica fallback, chiude/apre UDP e stampa diagnosi. \\
60--91 & Macchina a stati WiFi & Distingue connesso, transizione dopo perdita e tentativo di reconnessione ogni 3 s. \\
93--102 & Begin & Imposta station mode, evita persistenza NVS, abilita auto reconnect e avvia WiFi.begin. \\
104--144 & Invio & Verifica associazione, serializza V1, prova beginPacket, scrive, chiude e stampa una diagnostica periodica. \\
\bottomrule
\end{longtable}

\subsubsection{Funzioni e failure mode}
\codepath{isEmptyIp} riconosce 0.0.0.0. \codepath{configureUdpConnection} aggiorna \codepath{tankIP} e il socket. \par
\codepath{ensureWifiConnected} e' una macchina a stati: la prima connessione configura UDP, una perdita ferma il socket e ogni 3 s tenta reconnect/begin. \par
\codepath{UdpSender_send} tratta \codepath{beginPacket==0} come fallimento, ma non controlla \codepath{udp.begin}, \codepath{udp.print} o \codepath{udp.endPacket}; anche un invio allo stack locale non e' un ACK dal tank.

\subsubsection{Listing completo numerato}
\lstinputlisting[language=C++]{../controller-esp32/src/udpSender.cpp}

\subsubsection{Lettura riga per riga}
\begin{longtable}{p{0.10\textwidth}p{0.78\textwidth}}
\toprule
\textbf{Riga} & \textbf{Lettura riga per riga} \\
\midrule
\endhead
1 & Include il modulo locale udpSender.h e rende visibili la sua interfaccia e i suoi tipi. \\
3 & Importa l'API WiFi della board corrente. \\
4 & Importa il socket UDP fornito dal framework. \\
16 & Dichiara o inizializza una variabile: il tipo determina intervallo, durata e semantica del dato. \\
17 & Dichiara o inizializza una variabile: il tipo determina intervallo, durata e semantica del dato. \\
20 & Fissa la porta di comunicazione UDP\_\allowbreak{}PORT condivisa dai due firmware. \\
23 & Definisce una durata nominale in millisecondi per COORDINATE\_\allowbreak{}PRINT\_\allowbreak{}INTERVAL\_\allowbreak{}MS; viene rispettata solo quando il loop continua a girare. \\
26 & Definisce una durata nominale in millisecondi per WIFI\_\allowbreak{}RECONNECT\_\allowbreak{}INTERVAL\_\allowbreak{}MS; viene rispettata solo quando il loop continua a girare. \\
28 & Dichiara stato statico: dura per l'intera esecuzione e, se locale a funzione, conserva valore fra chiamate. \\
31 & Dichiara l'API tankIP senza implementarla: la definizione e' in un file .cpp o piu' avanti. \\
32 & Dichiara stato statico: dura per l'intera esecuzione e, se locale a funzione, conserva valore fra chiamate. \\
33 & Dichiara stato statico: dura per l'intera esecuzione e, se locale a funzione, conserva valore fra chiamate. \\
34 & Dichiara stato statico: dura per l'intera esecuzione e, se locale a funzione, conserva valore fra chiamate. \\
37 & Definisce la funzione isEmptyIp; i parametri costituiscono il suo input e il corpo seguente ne realizza il contratto. \\
38 & Termina la funzione e restituisce al chiamante il valore o il controllo indicato. \\
42 & Definisce la funzione configureUdpConnection; i parametri costituiscono il suo input e il corpo seguente ne realizza il contratto. \\
44 & Dichiara o inizializza una variabile: il tipo determina intervallo, durata e semantica del dato. \\
45 & Valuta una guardia condizionale prima di eseguire il blocco successivo. \\
46 & Aggiorna o inizializza uno stato; seguire la variabile permette di vedere l'effetto sul ciclo successivo. \\
49 & Invoca il socket UDP; un successo locale non equivale a conferma di ricezione remota. \\
50 & Invoca il socket UDP; un successo locale non equivale a conferma di ricezione remota. \\
51 & Legge o confronta il contatore temporale Arduino; la sottrazione unsigned e' rollover-safe ma richiede un loop che torni a eseguire. \\
53 & Emette diagnostica seriale; e' osservabilita' utile ma puo' aggiungere tempo al giro di loop. \\
54 & Emette diagnostica seriale; e' osservabilita' utile ma puo' aggiungere tempo al giro di loop. \\
55 & Emette diagnostica seriale; e' osservabilita' utile ma puo' aggiungere tempo al giro di loop. \\
56 & Emette diagnostica seriale; e' osservabilita' utile ma puo' aggiungere tempo al giro di loop. \\
57 & Emette diagnostica seriale; e' osservabilita' utile ma puo' aggiungere tempo al giro di loop. \\
62 & Definisce la funzione ensureWifiConnected; i parametri costituiscono il suo input e il corpo seguente ne realizza il contratto. \\
63 & Valuta una guardia condizionale prima di eseguire il blocco successivo. \\
64 & Valuta una guardia condizionale prima di eseguire il blocco successivo. \\
65 & Aggiorna o inizializza uno stato; seguire la variabile permette di vedere l'effetto sul ciclo successivo. \\
66 & Esegue l'istruzione o la chiamata indicata nel contesto della funzione corrente. \\
69 & Termina la funzione e restituisce al chiamante il valore o il controllo indicato. \\
72 & Valuta una guardia condizionale prima di eseguire il blocco successivo. \\
73 & Aggiorna o inizializza uno stato; seguire la variabile permette di vedere l'effetto sul ciclo successivo. \\
74 & Invoca il socket UDP; un successo locale non equivale a conferma di ricezione remota. \\
75 & Emette diagnostica seriale; e' osservabilita' utile ma puo' aggiungere tempo al giro di loop. \\
78 & Legge o confronta il contatore temporale Arduino; la sottrazione unsigned e' rollover-safe ma richiede un loop che torni a eseguire. \\
79 & Valuta una guardia condizionale prima di eseguire il blocco successivo. \\
80 & Aggiorna o inizializza uno stato; seguire la variabile permette di vedere l'effetto sul ciclo successivo. \\
81 & Emette diagnostica seriale; e' osservabilita' utile ma puo' aggiungere tempo al giro di loop. \\
85 & Valuta una guardia condizionale prima di eseguire il blocco successivo. \\
86 & Invoca l'API WiFi; associazione, gateway e stato rete sono gestiti dal framework, non dal protocollo applicativo. \\
90 & Termina la funzione e restituisce al chiamante il valore o il controllo indicato. \\
93 & Definisce la funzione UdpSender\_\allowbreak{}begin; i parametri costituiscono il suo input e il corpo seguente ne realizza il contratto. \\
94 & Emette diagnostica seriale; e' osservabilita' utile ma puo' aggiungere tempo al giro di loop. \\
95 & Emette diagnostica seriale; e' osservabilita' utile ma puo' aggiungere tempo al giro di loop. \\
97 & Invoca l'API WiFi; associazione, gateway e stato rete sono gestiti dal framework, non dal protocollo applicativo. \\
98 & Invoca l'API WiFi; associazione, gateway e stato rete sono gestiti dal framework, non dal protocollo applicativo. \\
99 & Invoca l'API WiFi; associazione, gateway e stato rete sono gestiti dal framework, non dal protocollo applicativo. \\
100 & Invoca l'API WiFi; associazione, gateway e stato rete sono gestiti dal framework, non dal protocollo applicativo. \\
101 & Legge o confronta il contatore temporale Arduino; la sottrazione unsigned e' rollover-safe ma richiede un loop che torni a eseguire. \\
104 & Definisce la funzione UdpSender\_\allowbreak{}send; i parametri costituiscono il suo input e il corpo seguente ne realizza il contratto. \\
105 & Valuta una guardia condizionale prima di eseguire il blocco successivo. \\
107 & Termina la funzione e restituisce al chiamante il valore o il controllo indicato. \\
110 & Dichiara o inizializza una variabile: il tipo determina intervallo, durata e semantica del dato. \\
113 & Serializza testo dentro un buffer con limite esplicito di dimensione. \\
114 & Esegue l'istruzione o la chiamata indicata nel contesto della funzione corrente. \\
116 & Valuta una guardia condizionale prima di eseguire il blocco successivo. \\
118 & Termina la funzione e restituisce al chiamante il valore o il controllo indicato. \\
121 & Invoca il socket UDP; un successo locale non equivale a conferma di ricezione remota. \\
122 & Invoca il socket UDP; un successo locale non equivale a conferma di ricezione remota. \\
125 & Valuta una guardia condizionale prima di eseguire il blocco successivo. \\
126 & Legge o confronta il contatore temporale Arduino; la sottrazione unsigned e' rollover-safe ma richiede un loop che torni a eseguire. \\
128 & Emette diagnostica seriale; e' osservabilita' utile ma puo' aggiungere tempo al giro di loop. \\
129 & Emette diagnostica seriale; e' osservabilita' utile ma puo' aggiungere tempo al giro di loop. \\
130 & Emette diagnostica seriale; e' osservabilita' utile ma puo' aggiungere tempo al giro di loop. \\
131 & Emette diagnostica seriale; e' osservabilita' utile ma puo' aggiungere tempo al giro di loop. \\
132 & Emette diagnostica seriale; e' osservabilita' utile ma puo' aggiungere tempo al giro di loop. \\
133 & Emette diagnostica seriale; e' osservabilita' utile ma puo' aggiungere tempo al giro di loop. \\
134 & Emette diagnostica seriale; e' osservabilita' utile ma puo' aggiungere tempo al giro di loop. \\
135 & Emette diagnostica seriale; e' osservabilita' utile ma puo' aggiungere tempo al giro di loop. \\
136 & Emette diagnostica seriale; e' osservabilita' utile ma puo' aggiungere tempo al giro di loop. \\
137 & Emette diagnostica seriale; e' osservabilita' utile ma puo' aggiungere tempo al giro di loop. \\
138 & Emette diagnostica seriale; e' osservabilita' utile ma puo' aggiungere tempo al giro di loop. \\
139 & Emette diagnostica seriale; e' osservabilita' utile ma puo' aggiungere tempo al giro di loop. \\
140 & Emette diagnostica seriale; e' osservabilita' utile ma puo' aggiungere tempo al giro di loop. \\
141 & Emette diagnostica seriale; e' osservabilita' utile ma puo' aggiungere tempo al giro di loop. \\
142 & Emette diagnostica seriale; e' osservabilita' utile ma puo' aggiungere tempo al giro di loop. \\
\bottomrule
\end{longtable}


\section{Riepilogo di studio e limiti pratici}

\begin{itemize}
\item Il firmware e' modulare: controller legge input, UDP trasporta, tank decide e moduli attuano.
\item Buffer e calcoli correnti sono generalmente limitati; i rischi principali non sono un overflow classico, ma fiducia in I/O, rete e stato fisico.
\item Non esistono interrupt applicativi: tutti i timer dipendono dal polling regolare.
\item Riarmo, debounce, fault AP/I2C e dead-time sono stati rafforzati nel software; il parser V1 e' invece mantenuto compatibile con la baseline e lascia A-08 aperto. Limiti torretta, timeout reale e impulso relay restano da provare elettricamente e meccanicamente.
\item Il PDF corrente unisce la baseline e il delta post-correzioni. Dopo il prossimo commit stabile conviene rigenerare completamente anche le tabelle riga per riga dei moduli modificati, per riallineare ogni numero di riga al listing dinamico.
\end{itemize}

\thispagestyle{fancy}

\end{document}
